% Standalone extraction of established MFP material. See SOURCE_MAP.json.
% Options for packages loaded elsewhere
% Options for packages loaded elsewhere
\PassOptionsToPackage{unicode}{hyperref}
\PassOptionsToPackage{hyphens}{url}
\PassOptionsToPackage{dvipsnames,svgnames,x11names}{xcolor}
%
\documentclass[
  11pt,
]{article}
\usepackage{xcolor}
\usepackage[a4paper,margin=25mm]{geometry}
\usepackage{amsmath,amssymb}
\setcounter{secnumdepth}{5}
\usepackage{iftex}
\ifPDFTeX
  \usepackage[T1]{fontenc}
  \usepackage[utf8]{inputenc}
  \usepackage{textcomp} % provide euro and other symbols
\else % if luatex or xetex
  \usepackage{unicode-math} % this also loads fontspec
  \defaultfontfeatures{Scale=MatchLowercase}
  \defaultfontfeatures[\rmfamily]{Ligatures=TeX,Scale=1}
\fi
\usepackage{lmodern}
\ifPDFTeX\else
  % xetex/luatex font selection
\fi
% Use upquote if available, for straight quotes in verbatim environments
\IfFileExists{upquote.sty}{\usepackage{upquote}}{}
\IfFileExists{microtype.sty}{% use microtype if available
  \usepackage[]{microtype}
  \UseMicrotypeSet[protrusion]{basicmath} % disable protrusion for tt fonts
}{}
\makeatletter
\@ifundefined{KOMAClassName}{% if non-KOMA class
  \IfFileExists{parskip.sty}{%
    \usepackage{parskip}
  }{% else
    \setlength{\parindent}{0pt}
    \setlength{\parskip}{6pt plus 2pt minus 1pt}}
}{% if KOMA class
  \KOMAoptions{parskip=half}}
\makeatother
% Make \paragraph and \subparagraph free-standing
\makeatletter
\ifx\paragraph\undefined\else
  \let\oldparagraph\paragraph
  \renewcommand{\paragraph}{
    \@ifstar
      \xxxParagraphStar
      \xxxParagraphNoStar
  }
  \newcommand{\xxxParagraphStar}[1]{\oldparagraph*{#1}\mbox{}}
  \newcommand{\xxxParagraphNoStar}[1]{\oldparagraph{#1}\mbox{}}
\fi
\ifx\subparagraph\undefined\else
  \let\oldsubparagraph\subparagraph
  \renewcommand{\subparagraph}{
    \@ifstar
      \xxxSubParagraphStar
      \xxxSubParagraphNoStar
  }
  \newcommand{\xxxSubParagraphStar}[1]{\oldsubparagraph*{#1}\mbox{}}
  \newcommand{\xxxSubParagraphNoStar}[1]{\oldsubparagraph{#1}\mbox{}}
\fi
\makeatother

\usepackage{color}
\usepackage{fancyvrb}
\newcommand{\VerbBar}{|}
\newcommand{\VERB}{\Verb[commandchars=\\\{\}]}
\DefineVerbatimEnvironment{Highlighting}{Verbatim}{commandchars=\\\{\}}
% Add ',fontsize=\small' for more characters per line
\usepackage{framed}
\definecolor{shadecolor}{RGB}{241,243,245}
\newenvironment{Shaded}{\begin{snugshade}}{\end{snugshade}}
\newcommand{\AlertTok}[1]{\textcolor[rgb]{0.68,0.00,0.00}{#1}}
\newcommand{\AnnotationTok}[1]{\textcolor[rgb]{0.37,0.37,0.37}{#1}}
\newcommand{\AttributeTok}[1]{\textcolor[rgb]{0.40,0.45,0.13}{#1}}
\newcommand{\BaseNTok}[1]{\textcolor[rgb]{0.68,0.00,0.00}{#1}}
\newcommand{\BuiltInTok}[1]{\textcolor[rgb]{0.00,0.23,0.31}{#1}}
\newcommand{\CharTok}[1]{\textcolor[rgb]{0.13,0.47,0.30}{#1}}
\newcommand{\CommentTok}[1]{\textcolor[rgb]{0.37,0.37,0.37}{#1}}
\newcommand{\CommentVarTok}[1]{\textcolor[rgb]{0.37,0.37,0.37}{\textit{#1}}}
\newcommand{\ConstantTok}[1]{\textcolor[rgb]{0.56,0.35,0.01}{#1}}
\newcommand{\ControlFlowTok}[1]{\textcolor[rgb]{0.00,0.23,0.31}{\textbf{#1}}}
\newcommand{\DataTypeTok}[1]{\textcolor[rgb]{0.68,0.00,0.00}{#1}}
\newcommand{\DecValTok}[1]{\textcolor[rgb]{0.68,0.00,0.00}{#1}}
\newcommand{\DocumentationTok}[1]{\textcolor[rgb]{0.37,0.37,0.37}{\textit{#1}}}
\newcommand{\ErrorTok}[1]{\textcolor[rgb]{0.68,0.00,0.00}{#1}}
\newcommand{\ExtensionTok}[1]{\textcolor[rgb]{0.00,0.23,0.31}{#1}}
\newcommand{\FloatTok}[1]{\textcolor[rgb]{0.68,0.00,0.00}{#1}}
\newcommand{\FunctionTok}[1]{\textcolor[rgb]{0.28,0.35,0.67}{#1}}
\newcommand{\ImportTok}[1]{\textcolor[rgb]{0.00,0.46,0.62}{#1}}
\newcommand{\InformationTok}[1]{\textcolor[rgb]{0.37,0.37,0.37}{#1}}
\newcommand{\KeywordTok}[1]{\textcolor[rgb]{0.00,0.23,0.31}{\textbf{#1}}}
\newcommand{\NormalTok}[1]{\textcolor[rgb]{0.00,0.23,0.31}{#1}}
\newcommand{\OperatorTok}[1]{\textcolor[rgb]{0.37,0.37,0.37}{#1}}
\newcommand{\OtherTok}[1]{\textcolor[rgb]{0.00,0.23,0.31}{#1}}
\newcommand{\PreprocessorTok}[1]{\textcolor[rgb]{0.68,0.00,0.00}{#1}}
\newcommand{\RegionMarkerTok}[1]{\textcolor[rgb]{0.00,0.23,0.31}{#1}}
\newcommand{\SpecialCharTok}[1]{\textcolor[rgb]{0.37,0.37,0.37}{#1}}
\newcommand{\SpecialStringTok}[1]{\textcolor[rgb]{0.13,0.47,0.30}{#1}}
\newcommand{\StringTok}[1]{\textcolor[rgb]{0.13,0.47,0.30}{#1}}
\newcommand{\VariableTok}[1]{\textcolor[rgb]{0.07,0.07,0.07}{#1}}
\newcommand{\VerbatimStringTok}[1]{\textcolor[rgb]{0.13,0.47,0.30}{#1}}
\newcommand{\WarningTok}[1]{\textcolor[rgb]{0.37,0.37,0.37}{\textit{#1}}}

\usepackage{longtable,booktabs,array}
\newcounter{none} % for unnumbered tables
\usepackage{calc} % for calculating minipage widths
% Correct order of tables after \paragraph or \subparagraph
\usepackage{etoolbox}
\makeatletter
\patchcmd\longtable{\par}{\if@noskipsec\mbox{}\fi\par}{}{}
\makeatother
% Allow footnotes in longtable head/foot
\IfFileExists{footnotehyper.sty}{\usepackage{footnotehyper}}{\usepackage{footnote}}
\makesavenoteenv{longtable}
\usepackage{graphicx}
\makeatletter
\newsavebox\pandoc@box
\newcommand*\pandocbounded[1]{% scales image to fit in text height/width
  \sbox\pandoc@box{#1}%
  \Gscale@div\@tempa{\textheight}{\dimexpr\ht\pandoc@box+\dp\pandoc@box\relax}%
  \Gscale@div\@tempb{\linewidth}{\wd\pandoc@box}%
  \ifdim\@tempb\p@<\@tempa\p@\let\@tempa\@tempb\fi% select the smaller of both
  \ifdim\@tempa\p@<\p@\scalebox{\@tempa}{\usebox\pandoc@box}%
  \else\usebox{\pandoc@box}%
  \fi%
}
% Set default figure placement to htbp
\def\fps@figure{htbp}
\makeatother





\setlength{\emergencystretch}{3em} % prevent overfull lines

\providecommand{\tightlist}{%
  \setlength{\itemsep}{0pt}\setlength{\parskip}{0pt}}




\usepackage[]{natbib}
\bibliographystyle{plainnat}


\usepackage{mathtools}
\usepackage{fvextra}
\RecustomVerbatimEnvironment{verbatim}{Verbatim}{breaklines=true,breakanywhere=true,fontsize=\small}
\RecustomVerbatimEnvironment{Highlighting}{Verbatim}{commandchars=\\\{\},breaklines=true,breakanywhere=true,fontsize=\small}
\allowdisplaybreaks[1]
\setlength{\emergencystretch}{3em}
\makeatletter
\@ifpackageloaded{caption}{}{\usepackage{caption}}
\AtBeginDocument{%
\ifdefined\contentsname
  \renewcommand*\contentsname{Table of contents}
\else
  \newcommand\contentsname{Table of contents}
\fi
\ifdefined\listfigurename
  \renewcommand*\listfigurename{List of Figures}
\else
  \newcommand\listfigurename{List of Figures}
\fi
\ifdefined\listtablename
  \renewcommand*\listtablename{List of Tables}
\else
  \newcommand\listtablename{List of Tables}
\fi
\ifdefined\figurename
  \renewcommand*\figurename{Figure}
\else
  \newcommand\figurename{Figure}
\fi
\ifdefined\tablename
  \renewcommand*\tablename{Table}
\else
  \newcommand\tablename{Table}
\fi
}
\@ifpackageloaded{float}{}{\usepackage{float}}
\floatstyle{ruled}
\@ifundefined{c@chapter}{\newfloat{codelisting}{h}{lop}}{\newfloat{codelisting}{h}{lop}[chapter]}
\floatname{codelisting}{Listing}
\newcommand*\listoflistings{\listof{codelisting}{List of Listings}}
\usepackage{amsthm}
\theoremstyle{definition}
\newtheorem{definition}{Definition}[section]
\theoremstyle{plain}
\newtheorem{theorem}{Theorem}[section]
\theoremstyle{plain}
\newtheorem{proposition}{Proposition}[section]
\theoremstyle{plain}
\newtheorem{lemma}{Lemma}[section]
\theoremstyle{remark}
\AtBeginDocument{\renewcommand*{\proofname}{Proof}}
\newtheorem*{remark}{Remark}
\newtheorem*{solution}{Solution}
\newtheorem{refremark}{Remark}[section]
\newtheorem{refsolution}{Solution}[section]
\makeatother
\makeatletter
\makeatother
\makeatletter
\@ifpackageloaded{caption}{}{\usepackage{caption}}
\@ifpackageloaded{subcaption}{}{\usepackage{subcaption}}
\makeatother
\usepackage{bookmark}
\IfFileExists{xurl.sty}{\usepackage{xurl}}{} % add URL line breaks if available
\urlstyle{same}
% fallback for those not using the hyperref driver hyperxmp:
\makeatletter
\@ifundefined{xmpquote}{\newcommand{\xmpquote}[1]{#1}}{}
\makeatother
\hypersetup{
  pdftitle={Mean-Field Peeling},
  pdfauthor={Amir Joudaki},
  colorlinks=true,
  linkcolor={blue},
  filecolor={Maroon},
  citecolor={Blue},
  urlcolor={Blue},
  pdfcreator={LaTeX via pandoc}}


\title{Mean-Field Peeling}
\usepackage{etoolbox}
\makeatletter
\providecommand{\subtitle}[1]{% add subtitle to \maketitle
  \apptocmd{\@title}{\par {\large #1 \par}}{}{}
}
\makeatother
\subtitle{A Gaussian Calculus for Fixed Finite Neural Derivative
Programs}
\author{Amir Joudaki}
\date{2026-10-11}
\begin{document}
\maketitle
\begin{abstract}
We present a constructive Gaussian calculus for fixed finite derivative
computations of Gaussian-initialized deep neural networks. A typed
language specifies coordinate operations, normalized contractions,
matrix and transpose reuse, represented updates, and finite derivatives
and flow jets. Exact finite-width differentiation compiles these
operations into primitive programs. A causal Gaussian source-response
recursion then identifies their infinite-width scalar limits, with
convergence in every finite \(L^p\) and uniform moment bounds. Singular
covariance is allowed without identifying distinct formal source
variables. A restricted initialization subclass admits a reduction to
moments of products of activation derivatives. We include the complete
proofs, worked reuse and moving-flow examples, symbolic geometry and
label-dependent feature-kernel jets, and a documented Python compiler
with its executable examples and tests. Depth, data dimensions,
derivative order and update count are fixed; no positive-time
reconstruction or growing-order guarantee is asserted.
\end{abstract}

\clearpage
\renewcommand*\contentsname{Table of contents}
{
\setcounter{tocdepth}{2}
\tableofcontents
\clearpage
}

\section{Introduction}\label{sec-introduction}

Mean-field peeling (MFP) turns a specified finite computation on a
Gaussian-initialized neural network into a finite graph of multivariate
Gaussian expectations. The computation may include coordinate
nonlinearities, products, sums, normalized averages, named matrix and
transpose calls, represented parameter updates, and the derivatives
admitted by the typed language below. Matrix identity and every reuse
are retained throughout the construction.

There are two mathematical operations. First, the product, chain and
adjoint rules express the requested network derivative exactly at each
finite width. Second, a Gaussian source-response recursion computes the
limit of that expression as width tends to infinity. For instance, the
reused-matrix observable with identity activation in
Equation~\ref{eq-mfp-reuse-square} has expectation \(2+1/n\) at width
\(n\), while the Gaussian graph evaluates to its limit \(2\).

The master statement is Theorem~\ref{thm-mfp-finite-derivative-program}.
Its program length, hidden depth, sample count, derivative order and
number of updates are fixed before the width limit. It gives convergence
of scalar observables in every finite \(L^p\), uniform moments and joint
empirical laws. The initialization normal form in
Proposition~\ref{prp-mfp-initialization-normal-form} has the additional
literal preactivation restriction stated there. The accompanying
compiler makes the construction inspectable and executable; its tests
check implementation identities rather than constituting a proof of the
theorem.

This paper collects the established MFP material of \emph{A Dynamical
Theory of Deep Learning}, together with the complete bounded-program
foundation used in its proof and the compiler guide. The raw-entry
moment mechanism is attributed to Golikov and Yang in
Section~\ref{sec-mfp-uniform-moments}. The presentation does not claim a
new general probabilistic program theorem. The accompanying source map
records the maintained book passages and the unchanged Python sources
used for this standalone edition.

The preliminary calculations are in Section~\ref{sec-preliminaries}. The
typed language, Gaussian evaluation, complete proof, derivative rules
and worked neural examples follow in
Section~\ref{sec-mfp-finite-derivative-programs}. The executable
interface is in Section~\ref{sec-compiler-guide}. Appendix
\ref{sec-docs-special-data-limits-l3785} supplies the bounded-derivative
foundation, including singular query Grams and causal scalar feedback.

Throughout, \(N(0,K)\) denotes the centered multivariate Gaussian law
with positive semidefinite covariance \(K\), including singular laws.
For finite vectors, pairings are divided by \(n\) when explicitly
displayed that way; the Frobenius matrix pairing is
\(\langle A,B\rangle_F=\operatorname{tr}(A^TB)\). The empirical law of a
same-type tuple \(v^1,\ldots,v^k\) is
\(n^{-1}\sum_i\delta_{(v^1_i,\ldots,v^k_i)}\). The distance
\(\mathcal W_2\) is the quadratic Wasserstein distance, that is, the
infimum of \((\mathbb E\|X-Y\|_2^2)^{1/2}\) over couplings of the two
laws.

\section{Gaussian conditioning and finite
calculations}\label{sec-preliminaries}

\subsection{One forward call followed by a
transpose}\label{sec-docs-gaussian-calculus-l19}

Let \(W=W^{(2)}(0)\) have independent \(\mathcal{N}(0,1/n)\) entries and
let \(1\) be the length-\(n\) vector of ones. Put \(y=W\mathbf{1}\). Its
entries are independent standard Gaussians. Take a continuous real
function \(g\) of polynomial growth and define

\[
q=W^T g(y),\qquad
a_n=\frac1n\sum_i y_i g(y_i),\qquad
\sigma_n^2=\frac1n\sum_i g(y_i)^2,\qquad
P=I-\frac{\mathbf1\mathbf1^T}{n}.
\]

Conditional on \(y\), the exact matrix law is

\begin{equation}\protect\phantomsection\label{eq-docs-gaussian-calculus-l34}{
W\ \overset{d}{=}\ \frac{y\mathbf1^T}{n}+\widetilde W P,
\qquad
q\ \overset{d}{=}\ a_n\mathbf1+\sigma_n P\gamma,
}\end{equation}

where \(\widetilde W\) has the original Gaussian matrix law and
\(\gamma\) is a standard Gaussian vector, each independent of \(y\) in
its corresponding representation. For the first equality, decompose each
Gaussian row into its component along \(1\) and its perpendicular
component. Their covariance vanishes, so they are independent; the first
component is determined by the observed row sum. Multiplying the
conditional matrix by \(g(y)\) on the left gives the second equality,
because \(\widetilde W^{\top}g(y)\) has conditional covariance
\(\sigma_n^2I\). In particular

\begin{equation}\protect\phantomsection\label{eq-docs-gaussian-calculus-l50}{
\mathbb E[q\mid y]=a_n\mathbf1,\qquad
\operatorname{Cov}(q\mid y)=\sigma_n^2P.
}\end{equation}

The random coefficient \(a_n\) is not replaced by its expectation at
finite width. If \(G\) is a standard scalar Gaussian, the elementary law
of large numbers gives

\begin{equation}\protect\phantomsection\label{eq-docs-gaussian-calculus-l60}{
a_n\longrightarrow a=\mathbb E[Gg(G)],\qquad
\sigma_n^2\longrightarrow\sigma^2=\mathbb E[g(G)^2]
\quad\text{in probability}.
}\end{equation}

All moments used here are finite by polynomial growth and Gaussian
tails. On the coupling in (\ref{eq-docs-gaussian-calculus-l34}), compare
\(q\) with \(a\mathbf{1}+\sigma\gamma\). Their RMS distance is at most

\begin{equation}\protect\phantomsection\label{eq-docs-gaussian-calculus-l71}{
|a_n-a|+|\sigma_n-\sigma|\frac{\|\gamma\|_2}{\sqrt n}
+\sigma_n\left|\frac1n\sum_i\gamma_i\right|,
}\end{equation}

which tends to zero in probability. This remains valid when
\(\sigma=0\). The empirical law of \(a+\sigma\gamma_i\) converges in
quadratic Wasserstein distance to the law of \(a+\sigma G\): weak
convergence follows by applying the law of large numbers to a countable
dense family of bounded Lipschitz tests, and the second moment converges
by the same law. On the real line, weak convergence together with
convergence of second moments yields quadratic Wasserstein convergence.
One direct verification truncates at a fixed radius, couples the bounded
laws using their quantiles, and lets the radius grow; convergence of
second moments makes the two squared-tail contributions vanish. Matching
corresponding coordinates and using
(\ref{eq-docs-gaussian-calculus-l71}) proves the same convergence for
the empirical law of \(q_i\).

If \(g\) is also continuously differentiable and \(g'\) has polynomial
growth, integration by parts in the Gaussian density gives

\begin{equation}\protect\phantomsection\label{eq-docs-gaussian-calculus-l92}{
a=\frac1{\sqrt{2\pi}}\int xg(x)e^{-x^2/2}dx
=\frac1{\sqrt{2\pi}}\int g'(x)e^{-x^2/2}dx.
}\end{equation}

The boundary term is zero because a polynomial times the Gaussian
density tends to zero. The transpose output therefore contains a
response \(\mathbb{E}[g'(G)]\) forced by the preceding forward call, in
addition to fresh Gaussian randomness. Treating the transpose as a new
independent matrix would erase that response. The two neuron populations
remain distinct; pairing row \(i\) with column \(i\) is not part of this
empirical-law statement.

\subsection{Exact conditioning after both directions have been
used}\label{sec-docs-gaussian-calculus-l105}

Suppose past calls have revealed \(WV=Y\) and \(W^{\top}U=R\). Columns
of \(V\) are input vectors in the first population; columns of \(U\) are
input vectors in the second. Assume both lists are linearly independent,
allowing either list to be empty. All matrix dimensions are finite.
Define

\[
P_V=V(V^TV)^{-1}V^T,\qquad P_U=U(U^TU)^{-1}U^T,
\]

\[
M=Y(V^TV)^{-1}V^T+U(U^TU)^{-1}R^T(I-P_V).
\]

Conditional on the complete transcript, the exact remaining matrix law
is

\begin{equation}\protect\phantomsection\label{eq-docs-gaussian-calculus-l121}{
W\ \overset d=\ M+(I-P_U)\widetilde W(I-P_V).
}\end{equation}

To verify the mean, use the compatibility identity
\(U^{\top}Y=R^{\top}V\), which holds because both lists were observed
from the same matrix. It gives \(MV=Y\) and \(M^{\top}U=R\). A
homogeneous perturbation preserves both observations exactly when
\(AV=0\) and \(A^{\top}U=0\), equivalently \(A=(I-P_U)A(I-P_V)\). The
displayed \(M\) is Frobenius-orthogonal to every such perturbation: its
first summand has right support in \(\operatorname{span}(V)\), and its
second has left support in \(\operatorname{span}(U)\). Thus it is the
minimum-norm point of the affine constraint set. Orthogonal components
of an isotropic finite Gaussian vector are independent. Applying this
fact to the vector of all matrix entries proves
(\ref{eq-docs-gaussian-calculus-l121}), including the scale \(1/n\) of
the remaining covariance.

This reasoning also holds for a finite adaptive transcript when every
new input is a measurable function of roots independent of \(W\) and of
earlier answers, and is chosen before its own answer is seen. Condition
on the roots first. Inductively, the next query is fixed by the past,
and conditioning its Gaussian answer adds a linear constraint to the
remaining Gaussian subspace. The orthogonal-decomposition argument
applies there. Derived variables measurable from these recorded roots
and answers add no further information. This induction does not allow an
unrecorded observation of \(W\) or an input chosen using its future
answer.

Dependent query columns can be discarded: each is an exactly known
linear combination of retained columns and its answer is the same
combination of retained answers. This does not justify taking limits of
ill-conditioned Gram inverses. Infinite-width arguments with nearly
dependent directions require an additional stability argument.

For a new input \(v\), (\ref{eq-docs-gaussian-calculus-l121}) explicitly
yields

\begin{equation}\protect\phantomsection\label{eq-docs-gaussian-calculus-l155}{
\begin{gathered}
Wv\ \overset d=\ Y\alpha+U\beta+
\frac{\|(I-P_V)v\|_2}{\sqrt n}(I-P_U)\gamma,
\\
\alpha=(V^TV)^{-1}V^Tv,\quad
\beta=(U^TU)^{-1}R^T(I-P_V)v.
\end{gathered}
}\end{equation}

The first response preserves previous forward answers, the second
restores the overlap demanded by previous transpose answers, and the
random part preserves both sets of constraints. Exchanging the two
populations gives the transpose formula. This is a finite exact
identity, not an iid assertion about coordinates after reuse.

\subsection{Derivatives along the gradient
direction}\label{sec-docs-gaussian-calculus-l170}

Let \(F:\mathbb{R}^p\to\mathbb{R}\) be \(C^3\) on an open set. At a
point write \(g=\nabla F\), \(H=\nabla^2F\) and \(T=\nabla^3F\). For a
differentiable scalar function \(u\), define the directional
differential operator \(\mathcal{D}_Fu=\langle g,\nabla u\rangle\).
Repeated application gives

\begin{equation}\protect\phantomsection\label{eq-docs-gaussian-calculus-l177}{
\mathcal D_F F=\|g\|_2^2,\qquad
\mathcal D_F^2F=2\langle g,Hg\rangle,\qquad
\mathcal D_F^3F=2T[g,g,g]+4\|Hg\|_2^2.
}\end{equation}

The first formula is the definition. Differentiating
\(\lVert g\rVert^2\) in direction \(g\) uses \(dg[g]=Hg\) and gives the
second formula. To differentiate \(\langle g,Hg\rangle\), differentiate
its three factors. The derivative of \(H\) contributes \(T[g,g,g]\). The
two derivatives of \(g\) contribute \(\langle Hg,Hg\rangle\) and
\(\langle g,H^2g\rangle\), both equal to \(\lVert Hg\rVert^2\) because
the Hessian is symmetric. Multiplication by two gives the last identity.

For a constant positive definite learning metric \(D\), change
coordinates to \(\theta=D^{1/2}q\) and apply
(\ref{eq-docs-gaussian-calculus-l177}) to \(F(D^{1/2}q)\). Its Euclidean
gradient flow in \(q\) is exactly the metric gradient flow in
\(\theta\). This supplies the metric version without confusing raw
gradients, mobilities and residuals.

If \(\theta(s)\) follows \(\theta'=\nabla F(\theta)\), the chain rule
identifies \(d^jF(\theta(s))/ds^j\) with \(\mathcal{D}_F^jF\) whenever
the stated derivatives exist. For squared loss the physical velocity has
the additional residual factor, so conversion to physical derivatives
also differentiates that factor. Finite-order identities alone provide
neither an analytic Taylor expansion nor permission to exchange width
limits and time differentiation.

\subsection{Exact polynomial Gaussian
expectations}\label{sec-docs-gaussian-calculus-l203}

Let \(\Sigma\) be a finite symmetric positive semidefinite matrix and
\(X\) a centered Gaussian vector with that covariance. For a vector
\(\alpha\) of nonnegative integer powers, put
\(M(\alpha)=\mathbb{E}[\prod_jX_j^{\alpha_j}]\). The zero-degree moment
is one; every odd-total-degree moment is zero by the symmetry \(X\) and
\(-X\). If \(\alpha_i>0\), Gaussian integration by parts gives

\begin{equation}\protect\phantomsection\label{eq-docs-gaussian-calculus-l211}{
M(\alpha)=\sum_{j=1}^p(\alpha_j-\mathbf1_{j=i})\Sigma_{ij}
 M(\alpha-e_i-e_j),
}\end{equation}

where terms with zero coefficient are omitted, so negative exponents
never occur. To prove the integration-by-parts identity even for
singular covariance, write \(X=AG\), where \(G\) is a standard Gaussian
vector and \(AA^{\top}=\Sigma\). Such a square root exists by
diagonalizing the finite symmetric matrix. For any polynomial \(P\),
integrate each component of \(G\) against its Gaussian density.
Polynomial growth makes all boundary terms zero, and the chain rule
gives

\[
\mathbb E[X_iP(X)]
=\sum_{k,j}A_{ik}A_{jk}\mathbb E[\partial_jP(X)]
=\sum_j\Sigma_{ij}\mathbb E[\partial_jP(X)].
\]

Apply this identity to \(P(X)=\prod_jX_j^{\alpha_j-\mathbf{1}_{j=i}}\)
to obtain (\ref{eq-docs-gaussian-calculus-l211}). Every nonzero summand
reduces total degree by two. Induction therefore proves termination and
correctness of the moment recurrence. With rational covariance all its
arithmetic is rational, and linearity evaluates every polynomial with
rational coefficients. The exact API \texttt{pde.gaussian\_moment}
implements this recurrence using integer and \(\mathrm{Fraction}\)
inputs; it does not approximate a nonpolynomial activation by a
polynomial or identify a training limit.

For completeness, its covariance validation also has an exact finite
proof. At a symmetric block matrix

\[
\begin{pmatrix}a&b^T\\ b&C\end{pmatrix},
\]

a negative \(a\) contradicts nonnegativity on the first basis vector. If
\(a=0\), positive semidefiniteness forces \(b=0\): otherwise the
quadratic form on \((t,v)\) with \(b^{\top}v\) nonzero changes sign for
a suitable large positive or negative \(t\). The remaining condition is
then exactly that \(C\) is positive semidefinite. If \(a>0\), completing
the square gives

\[
a t^2+2t b^Tv+v^TCv
=a\left(t+\frac{b^Tv}{a}\right)^2
+v^T\left(C-\frac{bb^T}{a}\right)v.
\]

Thus the condition is exactly positive semidefiniteness of the Schur
complement. Recursing proves the rational validator without a floating
tolerance, matrix inverse at a zero pivot, or a nonsingularity
assumption. This validates a finite covariance; it is not a proof about
the stability of adaptive query inverses.

\section{Fixed finite Gaussian derivative
programs}\label{sec-mfp-finite-derivative-programs}

The conditioning identities above become an operational calculus once
the permitted finite computations and derivatives are specified. The
order is to differentiate the finite network, expand every represented
matrix action, and then evaluate a finite Gaussian expectation graph.
The first two steps are exact at finite width. The last step gives a
width limit, including convergence of expectations. Its scope is a fixed
computation: depth, derivative order and update count do not grow with
width.

\subsection{Typed language and limit
theorem}\label{sec-mfp-program-language}

\begin{definition}[Finite Gaussian
program]\protect\hypertarget{def-mfp-primitive-program}{}\label{def-mfp-primitive-program}

A vector type labels one copy of \(\mathbb R^n\). Fix finitely many
types and named initialized matrices
\(W:\mathbb R^n_\tau\to\mathbb R^n_\sigma\), with \(\tau\ne\sigma\).
Each matrix has independent \(N(0,1/n)\) entries, and different named
matrices are independent. A transpose call uses that same matrix. For
each type fix a finite Gaussian root tuple whose coordinate tuples are
iid over neuron indices. Its mean and covariance are fixed; its
covariance may be singular. Root tuples of distinct types and all
matrices are independent. A constant vector means a vector with one
fixed scalar in every coordinate, not an arbitrary width-dependent
deterministic array.

A primitive program is a finite acyclic list of scalar and typed vector
instructions:

\begin{itemize}
\tightlist
\item
  \textbf{Coordinate map:} \(v_i=F(u_i^1,\ldots,u_i^k;s_1,\ldots,s_b)\),
  with all vectors of the same type.
\item
  \textbf{Normalized average:}
  \(s=n^{-1}\sum_i F(u_i^1,\ldots,u_i^k;t_1,\ldots,t_b)\).
\item
  \textbf{Scalar arithmetic:} sums and products of earlier scalars.
\item
  \textbf{Named matrix action:} \(Wh\) or \(W^Tu\), with the declared
  input and output types.
\end{itemize}

Initial scalars are deterministic. Every \(F\) is a finite expression
made from its arguments, fixed real constants, sums, products and
\(\phi^{(r)}\), \(r\ge0\), where

\begin{equation}\protect\phantomsection\label{eq-mfp-activation-growth}{
\begin{aligned}
\phi&\in C^\infty(\mathbb R),\\
|\phi^{(r)}(x)|&\le C_r(1+|x|)^{d_r}
\quad\text{for every }r\ge0.
\end{aligned}
}\end{equation}

All coefficients, root laws, expression lists and dimensions other than
\(n\) are fixed before \(n\to\infty\). There is no coordinate selection,
inverse operation, arbitrary tensor-index contraction or width-dependent
coefficient primitive.

\end{definition}

A current matrix and a matrix direction may be stored, respectively, as

\begin{equation}\protect\phantomsection\label{eq-mfp-represented-matrices}{
A=W_0+\sum_{\nu=1}^M c_\nu u_\nu v_\nu^T/n,
\qquad B=\sum_{\nu=1}^{M'} d_\nu p_\nu q_\nu^T/n,
}\end{equation}

with a fixed finite number of terms; the initialized part may be absent.
These are notation for primitive actions, not extra random-matrix laws:

\begin{equation}\protect\phantomsection\label{eq-mfp-rank-lowering}{
\begin{aligned}
Ah&=W_0h+\sum_\nu c_\nu u_\nu(v_\nu^Th/n),\\
A^Tu&=W_0^Tu+\sum_\nu c_\nu v_\nu(u_\nu^Tu/n).
\end{aligned}
}\end{equation}

Each pairing contracts only one type. The derivative language
additionally allows derivatives in finitely many deterministic inputs or
scalar parameters; fixed-order directional derivatives with program
vector directions and represented matrix directions; ambient gradients
with mobility \(n\) on trainable vector roots and mobility one on
trainable matrices, with fixed block multipliers; their admitted
contractions and directional derivatives; fixed numbers of gradient
updates; and finite gradient-flow jets. The exact rules are in
\hyperref[sec-mfp-physical-differentiation]{exact finite
differentiation}. No derivative of a sampling law is implicit in a
parameter derivative.

\begin{theorem}[Scalar limits and uniform
moments]\protect\hypertarget{thm-mfp-finite-derivative-program}{}\label{thm-mfp-finite-derivative-program}

Every derivative program just specified expands exactly at each finite
\(n\) to a finite primitive program. The construction in
\hyperref[sec-mfp-source-evaluation]{Gaussian source-response
evaluation} gives a deterministic scalar \(O_*\) for each scalar output
\(O_n\), and

\begin{equation}\protect\phantomsection\label{eq-mfp-scalar-lp-limit}{
\lim_{n\to\infty}\mathbb E|O_n-O_*|^p=0
\qquad(1\le p<\infty).
}\end{equation}

For every vector node \(v_n\), scalar node \(s_n\), and finite
\(p,q\ge1\),

\begin{equation}\protect\phantomsection\label{eq-mfp-uniform-node-moments}{
\sup_n\mathbb E\left(\frac1n\sum_i|v_{n,i}|^p\right)^q<\infty,
\qquad \sup_n\mathbb E|s_n|^q<\infty.
}\end{equation}

Every same-type tuple of vector nodes has a deterministic joint
empirical law converging in probability in \(\mathcal W_2\), for the
Euclidean metric on the tuple. Its law is the source-response law below.
These assertions hold jointly for finitely many programs sharing roots
and matrices, including singular query Grams.

The theorem asserts no infinite Taylor-series convergence, positive-time
reconstruction, common existence horizon for finite flows, growing-depth
or growing-update limit, or bound on compilation complexity beyond
finiteness. Vector coordinates do not converge to deterministic
constants.

\end{theorem}

The proof has three parts. A raw-entry differentiation estimate bounds
all moments without rank assumptions. The established bounded-derivative
Gaussian program law identifies each clipped program; the moment bounds
remove clipping uniformly in width. Exact finite differentiation closes
the language before this probabilistic argument is applied.

\subsection{Gaussian source-response
evaluation}\label{sec-mfp-source-evaluation}

Give each type \(\tau\) a probability space \(\Omega_\tau\) and
expectation \(\mathbb E_\tau\). Capital letters denote scalar
representatives of vector nodes. Roots have their declared laws;
coordinate maps act on these representatives and on previously computed
deterministic scalars. An average becomes the expectation on its type.

For a named matrix \(W:\tau\to\sigma\), list forward calls \(y_r=Wh_r\)
and reverse calls \(v_s=W^Tu_s\) in their program order. Their scalar
inputs are \(H_r\) on \(\Omega_\tau\) and \(U_s\) on \(\Omega_\sigma\).
Introduce centered Gaussian sources \(\xi_r\) on the output type
\(\sigma\), and \(\zeta_s\) on the output type \(\tau\), with

\begin{equation}\protect\phantomsection\label{eq-mfp-source-grams}{
\operatorname{Cov}(\xi_r,\xi_t)=\mathbb E_\tau[H_rH_t],\qquad
\operatorname{Cov}(\zeta_s,\zeta_v)=\mathbb E_\sigma[U_sU_v].
}\end{equation}

Different oriented matrix groups, including the two orientations of one
matrix, are independent Gaussian groups and are independent of the
roots. Return

\begin{equation}\protect\phantomsection\label{eq-mfp-source-response}{
\begin{aligned}
Y_r&=\xi_r+\sum_{s:\,\text{earlier reverse call}}
 U_s\,\mathbb E_\tau[\partial_{\zeta_s}H_r],\\
V_s&=\zeta_s+\sum_{r:\,\text{earlier forward call}}
 H_r\,\mathbb E_\sigma[\partial_{\xi_r}U_s].
\end{aligned}
}\end{equation}

The source covariances use uncentered input second moments. Thus an
input \(H=1\) has source variance one. The independence assertion
concerns sources; the output fields in
Equation~\ref{eq-mfp-source-response} generally depend on earlier groups
through their responses.

A formal source partial differentiates the explicit expression in the
named source arguments, holding all previously computed expectations,
covariance entries and scalar coefficients fixed. Paths through other
matrices' earlier calls remain part of that expression. An unavailable
source has derivative zero. Each call retains a distinct formal argument
even at singular covariance: compute partials before restricting to the
Gaussian support. No covariance inverse is an evaluation rule.

This is a causal construction. Every response integrand is already
defined, and the next covariance is an input-Gram extension, hence
positive semidefinite. A Gaussian law exists even when that extension is
singular. Finite sums, products, composition and differentiation
preserve polynomial growth of every derivative, by the finite product
and chain rules. Induction therefore bounds every field and its formal
derivatives by a polynomial in the finite root/source list, so all
required expectations exist. Collecting them gives a finite directed
acyclic graph

\begin{equation}\protect\phantomsection\label{eq-mfp-expectation-graph}{
\begin{aligned}
I_j&=\int F_j(g;I_1,\ldots,I_{j-1})\\
&\qquad N(0,K_j(I_1,\ldots,I_{j-1}))(dg),\\
O_*&=P(I_1,\ldots,I_N),
\end{aligned}
}\end{equation}

where root means are fixed shifts inside \(F_j\). The final \(P\) is a
polynomial because scalar outputs are assembled from averages, sums and
products. The dimensions, integrands and covariance entries are
specified by the program, rather than supplied as unknown neural laws.

\subsection{Uniform raw-entry moments}\label{sec-mfp-uniform-moments}

The following finite-entry cancellation argument is adapted from the
moment mechanism in Appendix J of \citet{golikov2022nongaussian}. Its
needed form is proved here; no external program theorem is used as a
hypothesis.

\begin{lemma}[Moments of values and raw-entry
derivatives]\protect\hypertarget{lem-mfp-raw-entry-moments}{}\label{lem-mfp-raw-entry-moments}

Temporarily allow all independent centered Gaussian matrix entries to
have different variances, each at most \(1/n\). For a fixed primitive
program with smooth maps whose every partial derivative has polynomial
growth, let \(\alpha\) be a list of \(k\) raw matrix-entry
differentiation variables, with repetition allowed. For each finite
\(p\ge1\) and \(k\ge0\),

\begin{equation}\protect\phantomsection\label{eq-mfp-raw-entry-bounds}{
\sup_{n,i,|\alpha|=k}\mathbb E|\partial^\alpha v_{n,i}|^p<\infty,
\qquad
\sup_{n,|\alpha|=k}\mathbb E|\partial^\alpha s_n|^p<\infty.
}\end{equation}

Constants are uniform over these variances and depend only on the fixed
graph, root laws and finitely many polynomial derivative bounds. They
are consequently uniform over any family of maps having common such
bounds at each order. Derivatives refer to the deterministic program as
a function of unconstrained raw entries, evaluated at the random array;
this includes entries of variance zero.

\end{lemma}

\protect\phantomsection\label{proof-mfp-raw-entry-moments}{}

\begin{proof}
Induct over instructions, proving all finite moments and derivative
orders simultaneously at each node. At fixed \(n\), values and
derivatives have polynomial growth in finitely many Gaussian variables,
so integrability is available before uniform bounds are proved. Roots
have Gaussian moments and positive-order matrix derivatives zero.
Repeated product and chain rules for coordinate and scalar maps give
finitely many products of earlier derivatives and polynomially bounded
factors. Hölder's inequality and higher moments of the earlier nodes
bound each term. This covers scalar feedback. For an average,
differentiate its finite sum and use

\begin{equation}\protect\phantomsection\label{eq-mfp-average-jensen}{
\mathbb E\left|\frac1n\sum_i a_i\right|^p
\le\frac1n\sum_i\mathbb E|a_i|^p.
}\end{equation}

For a matrix call \(y=Wv\), the product rule gives

\begin{equation}\protect\phantomsection\label{eq-mfp-matrix-derivative}{
\partial^\alpha y_i=\sum_jW_{ij}\partial^\alpha v_j+R_i.
}\end{equation}

The remainder has at most \(k\) terms, each an earlier-node derivative
of order \(k-1\): differentiating the linear matrix factor fixes its
column index, and a second derivative of that factor vanishes. Thus
\(R_i\) is bounded. A transpose call has the same argument with a column
in place of a row.

Write \(w_j=W_{ij}\) and \(z_j=\partial^\alpha v_j\). It suffices to
bound an even integer moment \(p\) of \(\sum_jw_jz_j\), since larger
even moments bound every finite moment. Expand it into ordered tuples
\((j_1,\ldots,j_p)\). For one tuple let \(b\) be its number of distinct
indices, let \(S\) be the indices that appear exactly once, and put
\(s=|S|\), \(M=\prod_{a=1}^p w_{j_a}\) and \(Z=\prod_{a=1}^p z_{j_a}\).
We will bound \(|\mathbb E[MZ]|\) by \(C n^{-b}\).

For \(u\in S\) define the commuting difference
\(\Delta_uZ=Z-Z|_{w_u=0}\). Centering and independence give
\(\mathbb E[MZ]=\mathbb E[M\prod_{u\in S}\Delta_u Z]\): in every other
term of the expanded product at least one singleton weight appears only
as its centered linear factor. Applying the fundamental theorem of
calculus once in each such entry yields

\begin{equation}\protect\phantomsection\label{eq-mfp-singleton-cancellation}{
\mathbb E[MZ]=\int_{[0,1]^s}
\mathbb E\left[M\prod_{u\in S}w_u\,
 (\partial_S Z)(W^{(t)})\right]dt.
}\end{equation}

Here \(W^{(t)}\) replaces each selected raw entry \(w_u\) by \(t_uw_u\)
in the entire array, consistently in every reuse and transpose. Empty
products and a zero-dimensional integral cover \(s=0\). Gaussian moments
give

\[
\mathbb E\left|M\prod_{u\in S}w_u\right|^2
\le C n^{-(p+s)}.
\]

The other factor in Equation~\ref{eq-mfp-singleton-cancellation} has a
bounded second moment, uniformly in \(n,t\) and the selected entries.
Its product-rule expansion contains only derivatives of the earlier node
\(v\), evaluated at independent centered Gaussian entries of variances
at most \(1/n\). The strengthened induction hypothesis applies.
Differentiate in raw coordinates first and then evaluate at \(W^{(t)}\);
these are not derivatives of the composed map \(W\mapsto Z(W^{(t)})\).
No factors \(t_u\) or divisions by \(t_u\) occur, even at zero.
Fixed-width polynomial integrability justifies the centering and
integration exchanges.

Cauchy--Schwarz bounds the tuple contribution by \(C n^{-(p+s)/2}\). The
multiplicities satisfy \(p\ge s+2(b-s)=2b-s\), giving the desired
\(C n^{-b}\). For a fixed equality pattern there are at most \(C_p n^b\)
tuples, and finitely many patterns for fixed \(p\). Summing proves a
uniform bound for the matrix-call moment and closes the induction. Every
bound uses finitely many earlier derivative bounds and Gaussian moments,
never a lower variance or a Gram inverse.
\end{proof}

Taking \(k=0\) proves scalar moment bounds, and Jensen gives

\[
\mathbb E\left(\frac1n\sum_i|v_i|^p\right)^q
\le\frac1n\sum_i\mathbb E|v_i|^{pq},
\]

which proves Equation~\ref{eq-mfp-uniform-node-moments}. We also use the
uniform operator-norm moments from
Equation~\ref{eq-docs-special-data-limits-l3859} and
Equation~\ref{eq-docs-special-data-limits-l3866}: for every fixed finite
\(a\), \(\sup_n\mathbb E\|W\|_{\mathrm{op}}^a<\infty\). Those estimates
follow from a \(1/4\)-net and Gaussian tails and apply to any fixed list
of initialized matrices and their transposes.

\subsection{Bounded maps, scalar feedback and removal of
clipping}\label{sec-mfp-clipping-proof}

\begin{lemma}[The bounded-derivative foundation on a finite typed
graph]\protect\hypertarget{lem-mfp-bounded-typed-extension}{}\label{lem-mfp-bounded-typed-extension}

For any fixed list of named matrices between distinct types, the source
rule Equation~\ref{eq-mfp-source-response} identifies empirical laws in
\(\mathcal W_2\) in probability for \(C^1\) coordinate maps with bounded
continuous first partials. Normalized averages of such maps, globally
Lipschitz scalar maps, and globally Lipschitz dependence on preceding
scalars are allowed. The roots here may have any iid
finite-second-moment tuple law of the stated independence types.

\end{lemma}

\protect\phantomsection\label{proof-mfp-bounded-typed-extension}{}

\begin{proof}
The deterministic-coefficient proof is
Theorem~\ref{thm-docs-special-data-limits-l3820}, with adaptive
conditioning in \hyperref[sec-docs-special-data-limits-l3884]{adaptive
Gaussian conditioning}, the source identity in
Lemma~\ref{lem-docs-special-data-limits-l3966} and singular-query
removal in Lemma~\ref{lem-docs-special-data-limits-l4011}. Its statement
lists adjacent-layer matrices. The following checks extend the proof to
the present finite typed graph without enlarging its analytic
hypotheses.

Condition on all revealed roots and the complete preceding transcript.
For each \emph{named} matrix retain only its own forward and reverse
constraints \(WV=Y\), \(W^TU=Q\). Initially the residual matrix laws are
independent Gaussians. A coordinate or scalar instruction is measurable
from the transcript. At the next call its input is fixed conditionally,
so its answer conditions only the queried residual matrix factor by one
linear observation; all other conditional factors are unchanged. This
proves by induction the conditional product structure used in
Equation~\ref{eq-docs-special-data-limits-l3896} for any list of names,
including parallel edges and type-graph cycles. Chronology is still
acyclic.

That conditioning formula and its new-query form
Equation~\ref{eq-docs-special-data-limits-l3916} use only the two
endpoint spaces and their earlier query columns. The finite-rank
projection error is \(\operatorname{rank}(U)/n\), regardless of other
edges. Thus the positive-Gram empirical-law induction, including its
conditional test variance and second-moment calculation, carries over
verbatim at each named call. The Gaussian integration-by-parts
cancellation in Lemma~\ref{lem-docs-special-data-limits-l3966} also uses
only this matrix's two lists. Dependencies through other matrices are
retained in the input's explicit source expression. This gives
Equation~\ref{eq-mfp-source-response} with separate oriented groups for
every name.

For singular Grams, add a distinct independent \(\varepsilon\)-Gaussian
input perturbation at each call, as in
Lemma~\ref{lem-docs-special-data-limits-l4011}. Each new squared
distance to its own preceding query span has limit at least
\(\varepsilon^2\). A finite union of the operator-bound events controls
all named matrices; the same finite instruction induction gives an
\(O(\varepsilon)\) RMS error. The proof's coupling by positive
semidefinite covariance square roots and its convergence of bounded
source derivatives do not use layer adjacency. It therefore removes
\(\varepsilon\) for this graph.

To include scalars, proceed causally, replacing each preceding scalar
value by its deterministic limit in an auxiliary program. A normalized
average of a bounded-derivative map is the normalized pairing of its
output with the constant-one vector, so the empirical \(\mathcal W_2\)
law determines that scalar limit. Globally Lipschitz scalar operations
preserve convergence. If \(s_n\to s_*\) and \(F(u;s)\) has a global
Lipschitz constant \(C\), then

\[
\begin{gathered}
\frac{\|F(u_n;s_n)-F(\bar u_n;s_*)\|_2}{\sqrt n}\\
{}\le C\left(\sum_j\frac{\|u_n^j-\bar u_n^j\|_2}{\sqrt n}\right.\\
\left.{}+\sum_b|s_{n,b}-s_{*,b}|\right).
\end{gathered}
\]

For averages, Jensen gives the same estimate for their absolute error.
Matrix errors are bounded by the initialized operator norm times the
input RMS error. This proves the scalar and vector comparison by finite
induction; equivalently, scalar arguments can be broadcast and then
inserted into an ordinary same-type coordinate map. It is the mechanism
of \hyperref[sec-docs-special-data-limits-l4054]{causal scalar feedback}
with this additional globally Lipschitz coordinate dependence checked
explicitly. The deterministic scalar construction is causal, not a
fixed-point assumption. In its source partials the substituted scalars
are held fixed. The same Lipschitz estimates also include scalars when
removing the singular-query perturbations. This proves the stated
extension.
\end{proof}

This lemma leaves intact the broader finite-second-moment roots and
\(C^1\) hypotheses of the bounded-derivative foundation in
Section~\ref{sec-docs-special-data-limits-l3785}. The smooth Gaussian
theorem adds all finite scalar moments under its stronger assumptions.
Its scope remains the fixed Gaussian root laws and derivative language
specified here.

\protect\phantomsection\label{proof-mfp-finite-derivative-program}{}

\begin{proof}
First consider a primitive program. Choose a smooth cutoff \(\chi\)
equal to one on the unit ball and zero outside the radius-two ball. For
every coordinate or scalar operation set \(F_B(x)=\chi(x/B)F(x)\),
\(B\ge1\), including scalar arguments in \(x\). Scalar sums and products
can be clipped by this rule too; constants are unchanged. Each \(F_B\)
has bounded derivatives and is globally Lipschitz. For each derivative
order its polynomial growth bound is uniform in \(B\): product-rule
factors \(B^{-j}D^j\chi(x/B)\) are uniformly bounded. Also \(F_B\to F\)
with all derivatives on compact sets. The raw-entry moment lemma
consequently holds uniformly in both \(n\) and \(B\) for the original
and clipped programs.

Couple them using identical roots and matrices. We claim for
corresponding nodes that

\begin{equation}\protect\phantomsection\label{eq-mfp-clipping-errors}{
\begin{aligned}
\lim_{B\to\infty}\sup_n\mathbb E\frac{\|v_n-v_n^B\|_2^2}{n}&=0,\\
\lim_{B\to\infty}\sup_n\mathbb E|s_n-s_n^B|^2&=0.
\end{aligned}
}\end{equation}

Here are the error estimates needed to induct over the list. Regard a
vector coordinate as a random variable on the original probability space
times the uniform distribution on its neuron index. All finite moments
are uniformly bounded on this probability space. If a difference
vanishes in \(L^2\), interpolation with a uniformly bounded higher
moment makes it vanish in every fixed finite \(L^p\), uniformly in
width. This also applies to scalar errors and scalar broadcasts.

Polynomial growth of first derivatives gives

\begin{equation}\protect\phantomsection\label{eq-mfp-polynomial-difference}{
|F(x)-F(y)|\le C(1+|x|+|y|)^d|x-y|.
}\end{equation}

Split \(F(x)-F_B(y)\) into \(F(x)-F(y)\) and \(F(y)-F_B(y)\). The first
tends to zero in \(L^2\) by Hölder, the input errors upgraded to
\(L^4\), and uniform moments of the polynomial factor. The second is
bounded by \(C(1+|y|)^d\mathbf1_{|y|>B}\). Its squared expectation
vanishes uniformly by any sufficiently high moment and a power tail
bound. This proves the coordinate, scalar and, by Jensen, average steps.

For a matrix call with input error \(e\), correlated with \(W\) or not,

\begin{equation}\protect\phantomsection\label{eq-mfp-correlated-matrix-error}{
\begin{aligned}
\mathbb E\frac{\|We\|_2^2}{n}
&\le (\mathbb E\|W\|_{\mathrm{op}}^4)^{1/2}
 \left[\mathbb E(\|e\|_2^2/n)^2\right]^{1/2}\\
&\le C\left[\mathbb E\frac1n\sum_i|e_i|^4\right]^{1/2}
\longrightarrow0.
\end{aligned}
}\end{equation}

The operator moments and the interpolation just proved justify the last
line. Transposes obey the identical bound, and roots have zero error.
This completes the induction proving
Equation~\ref{eq-mfp-clipping-errors}, including causal scalar feedback.

For fixed \(B\), Lemma~\ref{lem-mfp-bounded-typed-extension} identifies
the clipped program and gives \(O_n^B\to O_*^B\) in probability. The
clipped Gaussian evaluator converges to the unclipped evaluator. To
prove this, induct over its expectation graph. Earlier convergent
coefficients stay bounded. All fields and each finite set of source
derivatives have common polynomial bounds in the finite root/source
list, because the clipped maps have common derivative bounds of
polynomial growth. Couple each source prefix with a fixed standard
Gaussian vector using its positive semidefinite covariance square root.
Covariance convergence implies square-root convergence, including at
rank loss, by the argument in
Lemma~\ref{lem-docs-special-data-limits-l4011}. Local convergence of the
maps and their derivatives gives pointwise convergence on this coupling;
bounded square roots give a common polynomial Gaussian dominator.
Dominated convergence therefore passes input Grams, response
expectations and all fixed field moments to their unclipped values. This
closes the coefficient induction and proves \(O_*^B\to O_*\).

For any \(\epsilon>0\) and sufficiently large fixed \(B\),

\[
\Pr(|O_n-O_*|>\epsilon)
\le\Pr(|O_n-O_n^B|>\epsilon/3)
   +\Pr(|O_n^B-O_*^B|>\epsilon/3).
\]

Take the width limit and then \(B\to\infty\), using
Equation~\ref{eq-mfp-clipping-errors} and Markov's inequality. This
proves convergence in probability. For finite \(p\), choose \(q>p\). The
raw-entry moment bound controls \(\mathbb E|O_n-O_*|^q\) uniformly. The
part of its \(p\)th moment on \(|O_n-O_*|>R\) is at most \(CR^{p-q}\);
the part on the complement tends to zero by convergence in probability
and boundedness (split once more at an arbitrarily small threshold).
This proves Equation~\ref{eq-mfp-scalar-lp-limit}.

For a same-type vector tuple, the coordinatewise coupling bounds the
squared empirical \(\mathcal W_2\) distance by its mean squared error.
The clipped empirical laws converge by the bounded lemma, and their
limiting laws converge in \(\mathcal W_2\) by the Gaussian dominated
convergence just proved. The triangle inequality proves the asserted
unclipped empirical-law convergence in probability. Finite unions of
programs give the joint assertions. Finally,
Proposition~\ref{prp-mfp-derivative-closure} below expands every
admitted derivative to a primitive graph before clipping; thus this
proof applies to the actual derivative, never to a derivative of a
clipped network.
\end{proof}

\subsection{Exact finite differentiation and
training}\label{sec-mfp-physical-differentiation}

Physical differentiation acts on the realized finite-dimensional
program, then evaluates at sampled parameters. A single Fréchet
derivative evaluates its direction at the base point and holds that
direction fixed. With \(\delta u=Du[\delta\theta]\) and
\(\delta s=Ds[\delta\theta]\), its rules are

\begin{equation}\protect\phantomsection\label{eq-mfp-forward-ad}{
\begin{aligned}
\delta F(u;s)
&=\sum_j F_{u^j}(u;s)\odot\delta u^j\\
&\quad+\sum_b F_{s_b}(u;s)\,\delta s_b,\\
\delta\left[\frac1n\sum_iF(u_i;s)\right]
&=\frac1n\sum_i\bigl[\sum_jF_{u^j}(u_i;s)\delta u_i^j\\
&\qquad+\sum_bF_{s_b}(u_i;s)\delta s_b\bigr],\\
\delta(Ah)&=(\delta A)h+A\delta h,\\
\delta(A^Tu)&=(\delta A)^Tu+A^T\delta u,\\
\delta(cuv^T/n)&=(\delta c)uv^T/n\\
&\quad+c(\delta u)v^T/n\\
&\quad+cu(\delta v)^T/n.
\end{aligned}
}\end{equation}

Products of vectors in these expressions are coordinatewise. Scalar
arithmetic uses the usual sum and product rules. Unlike a formal source
partial, physical differentiation propagates through finite averages and
their scalar feedback.

For reverse differentiation of a scalar \(O\), store vector adjoints
\(b_u=n\,\partial O/\partial u\) and scalar adjoints
\(a_s=\partial O/\partial s\). Initialize \(a_O=1\) and the others at
zero; visit instructions in reverse order and add every contribution:

For a coordinate map \(v=F(u;s)\), add for each argument

\begin{equation}\protect\phantomsection\label{eq-mfp-coordinate-adjoint}{
b_{u^j}\mathrel{+}=b_v\odot F_{u^j},\qquad
a_{s_b}\mathrel{+}=\frac1n\sum_i b_{v,i}F_{s_b}(u_i;s).
}\end{equation}

For an average \(t=n^{-1}\sum_iF(u_i;s)\), add

\begin{equation}\protect\phantomsection\label{eq-mfp-average-adjoint}{
b_{u^j}\mathrel{+}=a_tF_{u^j},\qquad
a_{s_b}\mathrel{+}=\frac{a_t}{n}\sum_iF_{s_b}(u_i;s).
}\end{equation}

For a matrix action \(v=Au\), add

\begin{equation}\protect\phantomsection\label{eq-mfp-matrix-adjoint}{
b_u\mathrel{+}=A^Tb_v,\qquad \nabla_AO\mathrel{+}=b_vu^T/n.
}\end{equation}

For a transpose action \(v=A^Tu\), add

\begin{equation}\protect\phantomsection\label{eq-mfp-transpose-adjoint}{
b_u\mathrel{+}=Ab_v,\qquad \nabla_AO\mathrel{+}=ub_v^T/n.
}\end{equation}

For example, matching \(dO=b_v^Tdv/n\) to
\(\langle\nabla_AO,dA\rangle_F+b_u^Tdu/n\) proves both matrix rules,
where \(\langle B,C\rangle_F=\operatorname{tr}(B^TC)\). Accumulation
includes every repeated occurrence of a shared parameter.

Compute a current ambient-matrix gradient with \(A\) treated as the
parameter \emph{before} substituting
Equation~\ref{eq-mfp-represented-matrices}. A derivative with respect to
the factors of previous updates is a history pullback, a different
requested function. Vector mobility \(n\kappa\) gives loss velocity
\(-\kappa b_u\), and matrix mobility \(\kappa\) gives
\(-\kappa\nabla_A\mathcal L\). Their contractions stay in the language:

\begin{equation}\protect\phantomsection\label{eq-mfp-gradient-contractions}{
\begin{aligned}
n\,\nabla_uO^T\nabla_uS
&=\frac{(b_u^O)^Tb_u^S}{n},\\
\left\langle\frac{uv^T}{n},\frac{pq^T}{n}\right\rangle_F
&=\left(\frac{u^Tp}{n}\right)\left(\frac{v^Tq}{n}\right).
\end{aligned}
}\end{equation}

For physical-time jets write \(u_{[r]}=\frac1{r!}(d^ru/dt^r)(0)\). At a
fixed requested order \(R\), coefficient extraction in the truncated
polynomial ring gives

\begin{equation}\protect\phantomsection\label{eq-mfp-jet-algebra}{
\begin{aligned}
(uv)_{[r]}&=\sum_{a+b=r}u_{[a]}v_{[b]},\qquad
(Au)_{[r]}=\sum_{a+b=r}A_{[a]}u_{[b]},\\
\phi(u)_{[r]}&=[t^r]\sum_{k=0}^r\frac{\phi^{(k)}(u_{[0]})}{k!}
                 \left(\sum_{j=1}^r u_{[j]}t^j\right)^k.
\end{aligned}
}\end{equation}

The \(k=0\) term is the constant \(\phi(u_{[0]})\). Averages act
coefficientwise, and mixed finite jets use multivariate truncated
polynomials. For a smooth finite-dimensional field
\(\dot\theta=V(\theta)\), compute

\begin{equation}\protect\phantomsection\label{eq-mfp-flow-jet-recursion}{
(r+1)\theta_{[r+1]}=[t^r]V\left(\sum_{j=0}^r\theta_{[j]}t^j\right).
}\end{equation}

All state dependence in \(V\) participates. In particular
\(O''=D^2O[V,V]+DO[DV[V]]\); the first term alone describes the straight
line with frozen initial direction. This is the general version of
\hyperref[sec-docs-gaussian-calculus-l170]{derivatives along the
gradient direction}; the neural applications below use this same finite
recursion.

\begin{proposition}[Exact compilation of admitted
derivatives]\protect\hypertarget{prp-mfp-derivative-closure}{}\label{prp-mfp-derivative-closure}

The derivative constructions above, fixed gradient updates and
fixed-order flow jets expand to primitive programs, exactly at every
finite width.

\end{proposition}

\protect\phantomsection\label{proof-mfp-derivative-closure}{}

\begin{proof}
Every partial derivative of an allowed coordinate expression is again
such an expression by finite chain and product rules. Thus
Equation~\ref{eq-mfp-forward-ad} and the reverse table introduce only
allowed coordinate maps, averages and matrix actions. Differentiating a
represented matrix term leaves finitely many represented terms;
Equation~\ref{eq-mfp-rank-lowering} expands their actions, and
Equation~\ref{eq-mfp-gradient-contractions} expands their admitted
contractions. A vector update subtracts a program vector, while a matrix
update adds finitely many normalized rank-one terms. Induction over a
fixed number of simultaneous steps proves exact closure, with the
ambient loss gradient formed before each current-state substitution.

Finite Taylor's theorem, not analyticity, proves
Equation~\ref{eq-mfp-jet-algebra} to any prescribed finite order.
Comparing coefficients in \(\dot\theta=V(\theta)\) gives
Equation~\ref{eq-mfp-flow-jet-recursion}. A smooth finite-dimensional
field has a local solution: on a sufficiently small closed ball its
derivative is bounded, and its integral map on continuous paths in that
ball is a contraction for sufficiently short time. Repeated
differentiation yields the finite jets. No common width-independent
interval is used. Finite induction on jet order therefore proves the
remaining closure assertion. Applying the limit theorem only after
expansion avoids any interchange of differentiation and the width limit.
\end{proof}

Explicitly frozen direction data can be handled as independent seeds on
an enlarged parameter space, with values assigned from the base state
after differentiation. For example, freezing \(x\) inside
\(n^{-1}\sum_i x_i^2\) sets its partial derivative in the remaining
\(x\) parameter to zero; it is no longer the ordinary derivative of the
original observable. Gaussian evaluation retains the seed's value and
explicit source dependence. Freezing physical seeds never freezes a
formal source path.

For neural applications let \(L\) be the hidden depth, \(m\) the number
of samples, \(d\) the input dimension and \(n\) the hidden width. Fix
\(L,m,d\) and data \(x_a\in\mathbb R^d\), \(y_a\in\mathbb R\),
\(1\le a\le m\). The first weight matrix \(W^{(1)}\) has shape
\(n\times d\), hidden matrices \(W^{(\ell)}\) have shape \(n\times n\)
for \(2\le\ell\le L\), and the stored readout \(W^{(L+1)}\) has length
\(n\). The data Gram is \(G_{ab}=x_a^Tx_b/d\), without whitening or a
nonsingularity requirement. Use the canonical forward model

\begin{equation}\protect\phantomsection\label{eq-mfp-neural-forward}{
\begin{aligned}
z_a^{(1)}&=W^{(1)}x_a/\sqrt d,\\
z_a^{(\ell)}&=W^{(\ell)}h_a^{(\ell-1)},\\
h_a^{(\ell)}&=\phi(z_a^{(\ell)}),\\
f_{n,a}&=(W^{(L+1)})^Th_a^{(L)}/n.
\end{aligned}
}\end{equation}

Take independent first-weight entries \(N(0,1)\), hidden entries
\(N(0,1/n)\), and stored readout entries \(N(0,\tau^2)\) for a fixed
\(\tau\). The first-matrix columns and readout are vector roots. This is
an order-one readout when \(\tau>0\), explicitly distinct from the
default width-dependent small-readout law with stored entries
\(N(0,1/n^2)\). With \(r_{n,a}=f_{n,a}-y_a\), full mean loss
\(\mathcal L_n=m^{-1}\sum_a r_{n,a}^2\), and block mobilities
\(n\kappa_1,\kappa_2,\ldots,\kappa_L,n\kappa_{L+1}\), its physical
gradient field is admitted. All samples share the first-matrix columns;
one must not replace correlated sample preactivations by independent
trainable parameters. Fixed data derivatives also propagate through
those shared columns. Exact small-readout finite recurrences remain
valid, but their width-dependent initialization is not a consequence of
this fixed-root-law theorem.

\subsection{Initialization activation-moment
form}\label{sec-mfp-initialization-normal-form}

\begin{proposition}[Restricted activation-moment
reduction]\protect\hypertarget{prp-mfp-initialization-normal-form}{}\label{prp-mfp-initialization-normal-form}

Construct the original feedforward preactivations of the centered,
zero-bias initialization in Equation~\ref{eq-mfp-neural-forward} first.
Consider a compiled derivative program at initialization in which every
nonlinear node is literally \(\phi^{(r)}(z_a^{(\ell)}(0))\) for one of
these original nodes. All other operations are the allowed polynomial
arithmetic, matrix actions and normalized contractions. Its scalar limit
is a finite polynomial in fixed model/data constants and recursively
computed Gaussian moments of products of activation derivatives. This
syntactic restriction excludes nonlinear functions of derivative fields
and of updated preactivations.

\end{proposition}

\protect\phantomsection\label{proof-mfp-initialization-normal-form}{}

\begin{proof}
No reverse call precedes the original forward pass. Hence its
preactivations are centered Gaussian sources, or linear combinations of
first-layer Gaussian roots. Their covariances are determined in order by

\begin{equation}\protect\phantomsection\label{eq-mfp-forward-covariances}{
\begin{aligned}
\Sigma^{(1)}_{ab}&=G_{ab}=x_a^Tx_b/d,\\
\Sigma^{(\ell+1)}_{ab}
&=\mathbb E_{Z\sim N(0,\Sigma^{(\ell)})}\phi(Z_a)\phi(Z_b).
\end{aligned}
}\end{equation}

All original preactivations, auxiliary roots and later named sources
form a finite joint Gaussian family on each type. A preactivation
represented as a linear combination can be adjoined as a Gaussian
integration coordinate with its induced covariance, even if this creates
singularity. Source partials still use the original named arguments;
adjoining an integration coordinate does not change their definition.

Inductively every field is a finite sum of expressions
\(cP(Z)\prod_\nu\phi^{(r_\nu)}(Z_{a_\nu})\), with \(P\) polynomial,
activation arguments designated original preactivations, and \(c\)
polynomial in previously computed scalar values. Fixed means of
auxiliary roots can be absorbed into the polynomial; designated
activation arguments are centered preactivations by the subclass
assumption. The class is closed under arithmetic and formal source
derivatives. A matrix action appends one Gaussian source and linear
combinations of old fields with coefficients equal to expectations in
the same class; its covariance entries are also such expectations. This
proves the symbolic induction.

For any centered Gaussian vector \(Z\) of covariance \(K\) and smooth
polynomial-growth \(F\),

\begin{equation}\protect\phantomsection\label{eq-mfp-stein-reduction}{
\mathbb E[Z_iF(Z)]=\sum_jK_{ij}\mathbb E[\partial_jF(Z)].
}\end{equation}

Indeed write \(Z=B\gamma\), \(BB^T=K\), with independent standard
Gaussian coordinates in \(\gamma\), and integrate by parts in each
coordinate. Polynomial growth removes boundary terms and justifies
Fubini; no invertibility is needed. Apply this identity to one explicit
Gaussian factor in \(P\). A derivative hitting the remaining polynomial
reduces its degree, and a derivative hitting an activation raises its
derivative order while still removing the selected explicit Gaussian
factor. Induction on polynomial degree terminates in covariance
monomials times atoms

\[
\mathbb E_{Z\sim N(0,K_j)}\prod_\nu\phi^{(r_{j,\nu})}(Z_{a_{j,\nu}}),
\]

whose covariances are already determined by earlier steps. Substituting
these finite expressions recursively proves the claimed polynomial form.
For polynomial-only programs this reduces to
Equation~\ref{eq-docs-gaussian-calculus-l211}. The scalar \(L^p\)
theorem identifies the result with the expectation limit; the reduction
is not a separate asymptotic assumption.
\end{proof}

\subsection{Worked programs and neural
observables}\label{sec-mfp-worked-programs}

\textbf{Reuse and the finite-width correction.} Let \(e\) be the
constant-one input, \(y=We\), \(u=\phi(y)\), \(v=W^Tu\) and
\(J_n=v^Tv/n\). For a standard Gaussian \(\gamma\), put
\(q=\mathbb E\phi(\gamma)^2\) and \(c=\mathbb E\phi'(\gamma)\). The
forward representative is \(\xi\sim N(0,1)\); the reverse one is
\(\zeta+c\) with \(\zeta\sim N(0,q)\). Thus

\begin{equation}\protect\phantomsection\label{eq-mfp-reuse-square}{
\begin{aligned}
J_n&\longrightarrow q+c^2\quad\text{in every finite }L^p,\\
\mathbb E J_n&=q+c^2\\
&\quad+\frac{\mathbb E[\gamma^2\phi(\gamma)^2]-q-c^2}{n}.
\end{aligned}
}\end{equation}

For the finite identity, use the exact projection law
Equation~\ref{eq-docs-gaussian-calculus-l34}: conditionally on \(y\) the
two orthogonal components of \(v\) give
\(\mathbb E[J_n\mid y]=(n^{-1}y^T\phi(y))^2+
(1-1/n)n^{-1}\|\phi(y)\|_2^2\). Taking expectations and using
\(\mathbb E[\gamma\phi(\gamma)]=c\) proves the formula. For
\(\phi(x)=x\) it gives \(2+1/n\), whereas Gaussian compilation returns
\(2\). Physical differentiation must accumulate both occurrences of
\(W\):

\[
\nabla_WJ_n=\frac2n\left[
uv^T+\bigl((Wv)\odot\phi'(y)\bigr)e^T\right].
\]

This follows by substituting \(dv=(dW)^Tu+W^T(\phi'(y)\odot(dW)e)\) into
\(dJ_n=2v^Tdv/n\).

\textbf{One normalized update.} For the explicitly different scalar loss
\(\mathcal L_n=n^{-1}\sum_i\phi((We)_i)\), set \(p=\phi'(We)\). The
differential gives \(\nabla_W\mathcal L_n=pe^T/n\). One step with fixed
\(\eta\) and matrix mobility one gives \(A=W-\eta pe^T/n\). Keep \(p\)
at its pre-update value. Exact lowering yields
\(A^Tp=W^Tp-\eta e(p^Tp/n)\) and \(Ae=We-\eta p\). With
\(q=\mathbb E\phi'(\gamma)^2\) and \(c=\mathbb E\phi''(\gamma)\),

\[
\frac{\|A^Tp\|_2^2}{n}\longrightarrow q+(c-\eta q)^2.
\]

The reverse source has variance \(q\) and response \(c\), while the
saved pairing tends to \(q\); this proves the formula. For
\(\phi(x)=x^4/4\) the answer is \(15+(3-15\eta)^2\).

\textbf{Singular duplicate calls.} Distinct calls \(y_1=We\), \(y_2=We\)
retain sources of covariance
\(\left(\begin{smallmatrix}1&1\\1&1\end{smallmatrix}\right)\). For
\(u=\phi(y_1+y_2)\) the two formal partials are separately
\(\phi'(\xi_1+\xi_2)\). On the Gaussian support
\((\xi_1,\xi_2)=(\gamma,\gamma)\), so

\[
\frac{\|W^Tu\|_2^2}{n}\longrightarrow
\mathbb E\phi(2\gamma)^2+4(\mathbb E\phi'(2\gamma))^2.
\]

This also follows by applying Equation~\ref{eq-mfp-reuse-square} to
\(\phi(2x)\). For \(u=\phi(y_1)-\phi(y_2)\) the source variance is zero
and the two response coefficients cancel. Its transpose output remains
exactly zero. Identifying the formal arguments before differentiation
would discard the information that makes this cancellation explicit.

\textbf{Moving direction.} Let \(x\in\mathbb R^n\) be a trainable vector
state, not a dataset input, with iid standard Gaussian initialization.
For \(\mathcal L_n=n^{-1}\sum_i x_i^4/4\) and vector mobility \(n\),
\(\dot x=-x^3\). The observable \(O_n=n^{-1}\sum_i x_i^2\) satisfies

\[
D^2O_n[V,V]=\frac2n\sum_i x_i^6,\qquad
DO_n[DV[V]]=\frac6n\sum_i x_i^6.
\]

The expected initial frozen and moving second derivatives therefore tend
to \(30\) and \(120\). The moving coefficient \(O_{[2]}\) tends to
\(60\), because jets include \(1/2!\). Directly,
\(x_i(t)=x_i(0)(1+2t x_i(0)^2)^{-1/2}\) verifies these finite
identities.

\textbf{A neural parameter derivative.} Take two hidden layers, scalar
input \(x=1\), no biases, and independent
\(W_i^{(1)},W_i^{(3)}\sim N(0,1)\), \(W_{ij}^{(2)}\sim N(0,1/n)\). This
explicitly uses order-one readout. The forward pass is
Equation~\ref{eq-mfp-neural-forward}, with \(z^{(1)}=W^{(1)}\). Define
\(b^{(1)}=n\nabla_{W^{(1)}}f_n\) and
\(O_n=\|b^{(1)}\|_2^2/n=n\|\nabla_{W^{(1)}}f_n\|_2^2\). Compute in order

\[
\begin{aligned}
q_1&=\mathbb E\phi(\gamma)^2,\\
Z&\sim N(0,q_1),\\
s_1&=\mathbb E\phi'(\gamma)^2,\\
s_2&=\mathbb E\phi'(Z)^2,\\
q_2&=\mathbb E\phi(Z)^2.
\end{aligned}
\]

Then \(O_n\to s_1s_2\) in every finite \(L^p\). Indeed exact
differentiation gives \(\delta^{(2)}=W^{(3)}\odot\phi'(z^{(2)})\) and
\(b^{(1)}=\phi'(z^{(1)})\odot W^{(2)T}\delta^{(2)}\). The transpose
response is \(\mathbb E[A\phi''(Z)]=0\), since the centered readout
representative \(A\) is independent of \(Z\). Its reverse source has
variance \(s_2\) and is independent of the first-layer root; the squared
product consequently has expectation \(s_1s_2\). Likewise the full
metric gradient norm tends to

\[
n\|\nabla_{W^{(1)}}f_n\|_2^2+
\|\nabla_{W^{(2)}}f_n\|_F^2+
n\|\nabla_{W^{(3)}}f_n\|_2^2
\longrightarrow s_1s_2+q_1s_2+q_2,
\]

using \(\nabla_{W^{(2)}}f_n=\delta^{(2)}h^{(1)T}/n\) and
\(n\nabla_{W^{(3)}}f_n=h^{(2)}\).

\textbf{Symbolic geometry and feature-kernel jets.} Keep two hidden
layers and order-one readout, now with \(m=d=2\). Fix real
\(\alpha,\beta,\gamma,y_1,y_2\) and set

\[
x_1/\sqrt2=(\alpha,0)^T,\quad
x_2/\sqrt2=(\beta,\gamma)^T,\qquad
G=\begin{pmatrix}\alpha^2&\alpha\beta\\
\alpha\beta&\beta^2+\gamma^2\end{pmatrix}.
\]

Let \(U,V\) be the shared columns of \(W^{(1)}\), each with iid
\(N(0,1)\) entries, so \(z_1^{(1)}=\alpha U\) and
\(z_2^{(1)}=\beta U+\gamma V\). With the forward pass above, use the
half-mean loss and its physical clock

\begin{equation}\protect\phantomsection\label{eq-mfp-kernel-example-clock}{
\begin{gathered}
\mathcal L_n=\frac1{2m}\sum_a(f_{n,a}-y_a)^2,\\
(\dot U,\dot V,\dot W^{(2)},\dot W^{(3)})\\
=-\bigl(n\nabla_U\mathcal L_n,n\nabla_V\mathcal L_n,\\
\qquad\nabla_{W^{(2)}}\mathcal L_n,n\nabla_{W^{(3)}}\mathcal L_n\bigr).
\end{gathered}
}\end{equation}

Relative to the full mean loss, the velocity is halved. The
last-hidden-layer feature kernel is \(K_{n,ab}=h_a^{(2)T}h_b^{(2)}/n\);
it is only the readout block of the tangent kernel. For \(r=0,1,2\)
request

\begin{equation}\protect\phantomsection\label{eq-mfp-kernel-example-jets}{
J_r=\lim_{n\to\infty}\mathbb E\left[
\frac1{r!}\left.\frac{d^r}{dt^r}K_{n,12}(t)\right|_{t=0}\right].
}\end{equation}

Each coefficient exists by
Theorem~\ref{thm-mfp-finite-derivative-program}. For generic \(\phi\),
form the loss gradient in the shared parameters, use the moving
recursion Equation~\ref{eq-mfp-flow-jet-recursion} through order two,
and compile the three scalar coefficients by
Equation~\ref{eq-mfp-source-response}. All nonlinear arguments in these
differentiated initialization graphs are original preactivations, so
Proposition~\ref{prp-mfp-initialization-normal-form} applies. The zeroth
coefficient is \(\mathbb E[\phi(Z_1)\phi(Z_2)]\) with
\(Z\sim N(0,\Sigma^{(2)})\) from
Equation~\ref{eq-mfp-forward-covariances}. Higher generic formulas are
finite expectation graphs and can be regenerated by the maintained
example; their expansion is not needed to define or evaluate the
algorithm.

For the compact linear check \(\phi(z)=z\), the results are

\begin{equation}\protect\phantomsection\label{eq-mfp-linear-kernel-jets}{
\begin{aligned}
J_0&=G_{12},\\
J_1&=0,\\
J_2&=\frac94(G_{11}y_1+G_{12}y_2)
                  (G_{12}y_1+G_{22}y_2).
\end{aligned}
}\end{equation}

Here is a direct derivation of the coefficient. Write \(W=W^{(2)}\),
\(w=W^{(3)}\), \(v_b=W^{(1)}x_b/\sqrt d\), \(z_b=Wv_b\), \(B=WW^T\) and
\(s=\sum_c y_cv_c/m\). Put \(\lambda_b=\sum_cG_{bc}y_c/m\) and
\(\rho=\sum_{c,e}y_cG_{ce}y_e/m^2\). At initialization \(f_{n,b}\to0\)
in every finite \(L^p\) because the readout is centered and independent.
Terms containing these scalar predictions vanish in the expected finite
jet limits: after expansion their remaining normalized factors have all
moments by the theorem. Replacing residual values by \(-y_c\) in the
first velocity gives

\[
\begin{aligned}
\dot W&=ws^T/n,\\
\dot w&=Ws,\\
\dot v_b&=\lambda_bW^Tw,\\
\dot z_b&=w k_b+\lambda_bBw,\\
k_b&=s^Tv_b/n\to\lambda_b.
\end{aligned}
\]

In the second derivative, retain the change of the residuals as well.
Their limiting derivatives are deterministic by the scalar theorem, and
multiply first feature derivatives that are odd in the centered readout;
their expected pairing with \(z_b\) is zero. Replacing these scalar
residual derivatives by their limits is legitimate by Hölder and the
all-moment bounds. Thus they contribute zero to \(J_2\).

For the remaining terms differentiate the displayed velocity. Since
\(\dot s=\rho W^Tw\), one has
\(\dot k_b=\rho w^Tz_b/n+\lambda_b w^TWs/n\), which tends to zero, and
\(\dot B=(w(Ws)^T+(Ws)w^T)/n\). Conditioning on \(W,v_1,v_2\) and using
\(\mathbb E[ww^T]=I\) gives the following limiting pairings:

\[
\mathbb E\frac{\dot z_a^T\dot z_b}{n}\longrightarrow
5\lambda_a\lambda_b,\qquad
\mathbb E\frac{z_a^T\ddot z_b}{n}\longrightarrow
(1+1+2)\lambda_a\lambda_b=4\lambda_a\lambda_b.
\]

For clarity, the first coefficient is
\(\lim\mathbb E\operatorname{tr}(I+B)^2/n=1+2+2=5\). The three surviving
terms in the second are respectively \((Ws)k_b\),
\(\lambda_b(Ws)(w^Tw/n)\) and \(\lambda_bBWs\). The terms containing
\(w\dot k_b\) or \(w((Ws)^Tw/n)\) vanish by Cauchy--Schwarz and the
\(O(n^{-1})\) conditional second moments of the normalized readout
pairings. Gaussian pairings give \(\mathbb E\operatorname{tr}B/n=1\) and
\(\mathbb E\operatorname{tr}B^2/n=2+1/n\); covariance of \(v_a,s\)
conditional on \(W\) is \(\lambda_a I\), yielding the stated
cross-pairings. For example, in \(\operatorname{tr}B^2=
\sum_{i,j,k,l}W_{ik}W_{jk}W_{jl}W_{il}\) its three Gaussian pairings
contribute \(n,n,1\). Consequently \(K_{n,ab}''\) has limiting
expectation \((4+2\cdot5+4)\lambda_a\lambda_b=18\lambda_a\lambda_b\).
Dividing by \(2!\) and putting \(m=2\) proves
Equation~\ref{eq-mfp-linear-kernel-jets}. The first coefficient vanishes
by readout centering after the same residual-value replacement. These
finite initial coefficients do not reconstruct the feature kernel at
positive time.

\subsection{Symbolic implementation}\label{sec-mfp-symbolic-interface}

The maintained interface is explicitly imported from
\texttt{pde.mfp\_compiler}; it builds typed finite graphs, performs
physical differentiation, and returns
Equation~\ref{eq-mfp-expectation-graph}. For example:

\begin{Shaded}
\begin{Highlighting}[]
\ImportTok{from}\NormalTok{ pde.mfp\_compiler }\ImportTok{import}\NormalTok{ Program}

\NormalTok{p }\OperatorTok{=}\NormalTok{ Program()}
\NormalTok{lower, upper }\OperatorTok{=}\NormalTok{ p.vector\_type(}\StringTok{"lower"}\NormalTok{), p.vector\_type(}\StringTok{"upper"}\NormalTok{)}
\NormalTok{W }\OperatorTok{=}\NormalTok{ p.matrix(}\StringTok{"W"}\NormalTok{, lower, upper)}
\NormalTok{z }\OperatorTok{=}\NormalTok{ W }\OperatorTok{@}\NormalTok{ p.one(lower)}
\NormalTok{v }\OperatorTok{=}\NormalTok{ W.T }\OperatorTok{@}\NormalTok{ p.phi(z)}
\NormalTok{answer }\OperatorTok{=}\NormalTok{ p.}\BuiltInTok{compile}\NormalTok{(p.inner(v, v), preactivations}\OperatorTok{=}\NormalTok{[z])}
\BuiltInTok{print}\NormalTok{(answer.formula())}
\end{Highlighting}
\end{Shaded}

This returns the two activation moments in
Equation~\ref{eq-mfp-reuse-square}. \texttt{gradient} returns \(n\)
times vector gradients and ordinary matrix gradients;
\texttt{jets(...,\ moving=True)} returns factorial-normalized flow
coefficients. Supplying \texttt{preactivations} requests the restricted
form and checks its literal-node and Gaussian conditions. The caller
still supplies the architecture and activation-growth assumptions.

The compiler guide in Section~\ref{sec-compiler-guide} (also supplied as
\texttt{code/MFP\_CALCULUS.md}) gives executable reuse, derivative and
kernel-jet examples, supported operations and an independent
finite-array interpreter. The implementation uses exact rational
symbolic arithmetic; finite floats are converted by
\texttt{Fraction(str(value))}, and root covariances are concrete numeric
inputs. Symbolic geometry can instead enter through fixed scalar
coefficients such as \(\alpha,\beta,\gamma\) above. Generic integrals
remain explicit; the compiler performs no numerical quadrature. Higher
derivatives and expression expansion can grow rapidly. The theorem
proves the limit under its hypotheses; implementation tests check finite
identities and the emitted formulas.

\section{The symbolic compiler: interface and
examples}\label{sec-compiler-guide}

\texttt{pde.mfp\_compiler} builds a finite typed program, differentiates
it at finite width, and compiles its scalar output to a Gaussian
expectation graph. Import it explicitly; its names are not added to
\texttt{pde}'s top-level exports. The
\hyperref[sec-mfp-finite-derivative-programs]{finite-program theorem}
supplies the width-limit justification under its stated assumptions.
Software tests check the implementation, not that theorem.

\textbf{Physical differentiation is exact at finite width. Gaussian
compilation returns the infinite-width limit.} For example, the
identity-activation reuse observable below has finite expectation
\texttt{2\ +\ 1/n}, whereas \texttt{compile} returns \texttt{2}. No
numerical sampling or quadrature is performed. General activation
integrals remain explicit integrals.

Use Python 3.10+ and the maintained package's NumPy dependency. These
three new modules themselves use only the standard library; importing
their parent \texttt{pde} also imports its existing NumPy finite-network
API. From the \texttt{paper\_MFP/} root:

\begin{Shaded}
\begin{Highlighting}[]
\VariableTok{PYTHONPATH}\OperatorTok{=}\NormalTok{code }\ExtensionTok{python} \AttributeTok{{-}B}\NormalTok{ code/scripts/example\_mfp\_calculus.py}
\VariableTok{PYTHONPATH}\OperatorTok{=}\NormalTok{code }\ExtensionTok{python} \AttributeTok{{-}B}\NormalTok{ code/scripts/example\_mfp\_mlp\_derivative.py}
\VariableTok{PYTHONPATH}\OperatorTok{=}\NormalTok{code }\ExtensionTok{python} \AttributeTok{{-}B}\NormalTok{ code/scripts/example\_mfp\_kernel\_jets.py }\AttributeTok{{-}{-}activation}\NormalTok{ identity}
\VariableTok{PYTHONPATH}\OperatorTok{=}\NormalTok{code }\ExtensionTok{python} \AttributeTok{{-}B} \AttributeTok{{-}m}\NormalTok{ unittest discover }\AttributeTok{{-}s}\NormalTok{ code/tests }\AttributeTok{{-}p} \StringTok{\textquotesingle{}test\_mfp\_*.py\textquotesingle{}} \AttributeTok{{-}v}
\end{Highlighting}
\end{Shaded}

The scripts print text by default and accept \texttt{-\/-json}. To
regenerate all three kernel-jet variants into a fresh, explicitly chosen
directory:

\begin{Shaded}
\begin{Highlighting}[]
\VariableTok{PYTHONPATH}\OperatorTok{=}\NormalTok{code }\ExtensionTok{python} \AttributeTok{{-}B}\NormalTok{ code/scripts/example\_mfp\_kernel\_jets.py }\AttributeTok{{-}{-}activation}\NormalTok{ all }\AttributeTok{{-}{-}output{-}dir}\NormalTok{ build/kernel\_jets\_01}
\end{Highlighting}
\end{Shaded}

Create its parent first with \texttt{mkdir\ -p\ build}; the output
directory itself must not exist. It receives complete text formulas,
JSON graphs, and a manifest with source/output hashes and
Python/platform information. The producer reads no historical output,
study files or Git metadata. The default jet order is two; larger fixed
orders are permitted but can grow very expensive.

\subsection{Construct a reusable-transpose
observable}\label{construct-a-reusable-transpose-observable}

A named matrix has iid Gaussian entries of variance \texttt{1/n}. Calls
to \texttt{W.T} use the same parameter as calls to \texttt{W}. Both
types have width \texttt{n}.

\begin{Shaded}
\begin{Highlighting}[]
\ImportTok{from}\NormalTok{ pde.mfp\_compiler }\ImportTok{import}\NormalTok{ Program}

\NormalTok{p }\OperatorTok{=}\NormalTok{ Program()}
\NormalTok{lower, upper }\OperatorTok{=}\NormalTok{ p.vector\_type(}\StringTok{"lower"}\NormalTok{), p.vector\_type(}\StringTok{"upper"}\NormalTok{)}
\NormalTok{W }\OperatorTok{=}\NormalTok{ p.matrix(}\StringTok{"W"}\NormalTok{, lower, upper)}
\NormalTok{z }\OperatorTok{=}\NormalTok{ W }\OperatorTok{@}\NormalTok{ p.one(lower)}
\NormalTok{v }\OperatorTok{=}\NormalTok{ W.T }\OperatorTok{@}\NormalTok{ p.phi(z)}
\NormalTok{dag }\OperatorTok{=}\NormalTok{ p.}\BuiltInTok{compile}\NormalTok{(p.inner(v, v), preactivations}\OperatorTok{=}\NormalTok{[z])}
\BuiltInTok{print}\NormalTok{(dag.formula())}
\end{Highlighting}
\end{Shaded}

For a standard Gaussian scalar \texttt{Z}, the returned limit is
\texttt{E{[}phi(Z)**2{]}\ +\ E{[}phi\textquotesingle{}(Z){]}**2}. The
second term is the response caused by matrix reuse. With a polynomial
coordinate expression such as \texttt{v\ =\ W.T\ @\ z**3}, all Gaussian
moments reduce exactly and the output is \texttt{24}.

\texttt{p.phi(u,\ order=r)} is the actual derivative
\texttt{phi\^{}(r)(u)} of one named activation. Assume \texttt{phi} is
smooth and every derivative has polynomial growth. This analytic
condition is the caller's obligation, not something inferred from a name
or callback. Arbitrary admitted compositions can be compiled without
\texttt{preactivations}. Specifying that argument requests the
restricted initialization activation-moment form: each reachable
\texttt{phi} argument must be one of those literal nodes, and every
designated reachable representative must be a centered Gaussian linear
expression. The compiler checks these conditions; the caller identifies
the original initialization preactivations. It does not infer the
network architecture. A shifted or nonlinear designated argument is
rejected. Linear combinations receive auxiliary integration coordinates;
formal source derivatives still use the original named sources.

\subsection{Typed operations and input
conventions}\label{typed-operations-and-input-conventions}

\begin{description}
\tightlist
\item[\texttt{p.vector\_type(name)}]
A named population of width \texttt{n}
\item[\texttt{p.root(name,\ kind,\ variance=1,\ mean=0)}]
An iid Gaussian vector independent of other root declarations
\item[\texttt{p.roots(kind,\ names,\ covariance,\ means=None)}]
One jointly Gaussian tuple per neuron, allowing singular covariance
\item[\texttt{p.parameter(name)}]
A fixed deterministic real scalar, available to physical AD
\item[\texttt{p.broadcast(s,\ kind)}]
The constant-coordinate vector with value \texttt{s}
\item[\texttt{u+v}, \texttt{u*v}, \texttt{u**r}]
Addition, coordinate product and nonnegative integer power; scalars
broadcast
\item[\texttt{p.phi(u,\ r)}]
A coordinatewise activation derivative
\item[\texttt{p.mean(u)}, \texttt{p.inner(u,v)}]
\texttt{sum(u\_i)/n} and \texttt{u.T\ v/n}; pairings require the same
type
\item[\texttt{W\ @\ u}, \texttt{W.T\ @\ v}]
Both orientations of the same named matrix
\item[\texttt{p.rank\_one(u,v,c)}]
\texttt{c*u*v.T/n}, from the type of \texttt{v} to that of \texttt{u}
\item[\texttt{W\ +\ p.rank\_one(...)}]
One initialized base plus a finite represented update sum
\item[\texttt{p.frobenius(A,B)}]
Exact Frobenius product of two pure represented matrix directions
\end{description}

All vector types have equal width. Named matrices connect distinct
types. Retain the original object for shared parameters: duplicate
parameter names are rejected. Separate root declarations are
independent, including separate calls of \texttt{roots}; use a single
call for correlations within one type.

Constants accept Python integers, \texttt{Fraction}, and finite floats.
A float is interpreted as the exact decimal rational
\texttt{Fraction(str(value))}. All later symbolic constant arithmetic is
rational; arithmetic already evaluated by Python before entering the
builder is outside that guarantee. Root means and covariance entries are
concrete numbers. Covariance symmetry and positive semidefiniteness are
checked exactly, including singular laws and zero pivots.

The frontend does not accept symbolic covariance matrices. Symbolic
fixed geometry can instead be represented by explicit scalar
coefficients multiplying independent roots, as in the kernel example
below. This distinction preserves the shared trainable parameters and
their derivatives. Arbitrary fixed real coefficients can be scalar
parameter nodes; admissibility always requires them to remain fixed as
width grows.

Build coordinate expressions from these operations. For instance,
\texttt{p.phi(u*v+s,\ 2)*w} is a same-type expression when
\texttt{u,v,w} share a type and \texttt{s} is scalar. Nonlinear scalar
feedback can be written \texttt{p.mean(p.phi(p.broadcast(s,\ kind)))}.
Its limit may be a zero-dimensional Gaussian expectation, so the final
assembly remains polynomial in the returned expectation nodes and
deterministic parameters.

Cross-type products, array or text formulas, nonconstant division,
negative or fractional powers, and dense matrix directions are
unsupported. There is no coordinate selection, inverse, arbitrary
parameter-tensor contraction, or width-dependent normalization
primitive. The implementation supports arbitrary fixed finite program
sizes, derivative orders and update counts in principle; Python
recursion, memory use, rational bit sizes and expression expansion can
stop much smaller practical cases. There is no uniform-in-depth or
efficient high-order guarantee. A parameter named \texttt{n} does not
make a width-dependent program admissible. Fixed-order jets do not
reconstruct a positive-time trajectory or establish convergence of an
infinite Taylor series.

\subsection{Physical derivatives and frozen seed
data}\label{physical-derivatives-and-frozen-seed-data}

\texttt{p.directional(O,\ directions)} takes one derivative of the
finite primitive graph. Dictionary keys are vector-root nodes, scalar
parameter nodes, or named matrix parameters. Node directions must match
their parameter type; matrix directions must be pure represented rank
sums. Inserted directions are held fixed for that one derivative.
Unlisted parameters have zero direction.

\texttt{p.derivatives(O,\ directions,\ order,\ moving=False)} returns
derivatives indexed \texttt{0,...,order}. With \texttt{moving=False},
directions are fixed across all repeated derivatives, giving the
straight path through the initial point. With \texttt{moving=True},
repeated differentiation includes the state dependence of the vector
field. Its second derivative is
\texttt{D²O{[}V,V{]}\ +\ DO{[}DV{[}V{]}{]}}.

\begin{Shaded}
\begin{Highlighting}[]
\ImportTok{from}\NormalTok{ pde.mfp\_compiler }\ImportTok{import}\NormalTok{ Program}

\NormalTok{p }\OperatorTok{=}\NormalTok{ Program()}
\NormalTok{kind }\OperatorTok{=}\NormalTok{ p.vector\_type(}\StringTok{"neurons"}\NormalTok{)}
\NormalTok{x }\OperatorTok{=}\NormalTok{ p.root(}\StringTok{"x"}\NormalTok{, kind)}
\NormalTok{O, V }\OperatorTok{=}\NormalTok{ p.mean(x}\OperatorTok{**}\DecValTok{2}\NormalTok{), }\OperatorTok{{-}}\NormalTok{(x}\OperatorTok{**}\DecValTok{3}\NormalTok{)}
\NormalTok{frozen }\OperatorTok{=}\NormalTok{ p.derivatives(O, \{x: V\}, order}\OperatorTok{=}\DecValTok{2}\NormalTok{)}
\NormalTok{moving }\OperatorTok{=}\NormalTok{ p.derivatives(O, \{x: V\}, order}\OperatorTok{=}\DecValTok{2}\NormalTok{, moving}\OperatorTok{=}\VariableTok{True}\NormalTok{)}
\ControlFlowTok{assert} \BuiltInTok{str}\NormalTok{(p.}\BuiltInTok{compile}\NormalTok{(frozen[}\DecValTok{2}\NormalTok{]).output) }\OperatorTok{==} \StringTok{"30"}
\ControlFlowTok{assert} \BuiltInTok{str}\NormalTok{(p.}\BuiltInTok{compile}\NormalTok{(moving[}\DecValTok{2}\NormalTok{]).output) }\OperatorTok{==} \StringTok{"120"}
\ControlFlowTok{assert} \BuiltInTok{str}\NormalTok{(p.}\BuiltInTok{compile}\NormalTok{(p.jets(O, \{x: V\}, }\DecValTok{2}\NormalTok{, moving}\OperatorTok{=}\VariableTok{True}\NormalTok{)[}\DecValTok{2}\NormalTok{]).output) }\OperatorTok{==} \StringTok{"60"}
\end{Highlighting}
\end{Shaded}

Here \texttt{x} is the evolving vector state, not a data input. For the
scalar ODE \texttt{x\_dot=-x**3}, the coordinate solution is
\texttt{a/sqrt(1+2*t*a**2)}, so the second derivative of its square at
zero is \texttt{8*a**6}. A standard Gaussian sixth moment is 15. The
straight path \texttt{a-t*a**3} instead gives \texttt{2*a**6}. This
independently explains the values 120 and 30. \texttt{jets} divides
derivative \texttt{r} by \texttt{r!}.

\texttt{p.freeze(value)} declares an independent seed datum with the
base state's numerical value. Physical AD holds it fixed, while finite
evaluation and Gaussian compilation retain that value and its source
dependence. Thus \texttt{p.mean(p.freeze(x)**2)} has the same value as
\texttt{p.mean(x**2)}, but zero physical derivative with respect to
\texttt{x} under the declared convention. This is a partial derivative
on the enlarged space of parameters and seed data, evaluated at the
specified seeds. Inserting \texttt{freeze} into a primal network changes
the derivative being requested. Gaussian source partials still
differentiate the frozen expression's explicit source dependence.
Independence here concerns physical derivative coordinates, not
statistical independence of initialization values.

State substitution preserves the original seed expression:
\texttt{p.at(seed,\ state)} does not refresh it, and repeated gradient
steps keep that same seed. For \texttt{seed=p.freeze(x)} and
\texttt{loss=p.inner(x,seed)}, vector mobility \texttt{n} gives
\texttt{x\_after=x-steps*eta*seed}. To take a new snapshot of an
ordinary expression \texttt{value} at an updated state, explicitly use
\texttt{p.freeze(p.at(value,state))}. Freezing a represented matrix
direction means freezing its scalar coefficient and both vector factors;
their original dependencies are preserved in the same way. Finite
evaluation still obtains seed values from the supplied initial arrays,
rather than storing a numerical copy when the symbolic node is built.

For a general finite path, \texttt{curve\_jets} accepts already
factorial-normalized parameter coefficients. This complete example uses
\texttt{x(t)=x+t*x**2-t**2*x}, with both coefficients evaluated at the
base state and fixed on the path:

\begin{Shaded}
\begin{Highlighting}[]
\ImportTok{from}\NormalTok{ pde.mfp\_compiler }\ImportTok{import}\NormalTok{ Program}

\NormalTok{p }\OperatorTok{=}\NormalTok{ Program()}
\NormalTok{x }\OperatorTok{=}\NormalTok{ p.root(}\StringTok{"x"}\NormalTok{, p.vector\_type(}\StringTok{"neurons"}\NormalTok{))}
\NormalTok{jet }\OperatorTok{=}\NormalTok{ p.curve\_jets(p.mean(x}\OperatorTok{**}\DecValTok{2}\NormalTok{), \{x: (x}\OperatorTok{**}\DecValTok{2}\NormalTok{, }\OperatorTok{{-}}\NormalTok{x)\}, order}\OperatorTok{=}\DecValTok{2}\NormalTok{)}
\ControlFlowTok{assert} \BuiltInTok{str}\NormalTok{(p.}\BuiltInTok{compile}\NormalTok{(jet[}\DecValTok{2}\NormalTok{]).output) }\OperatorTok{==} \StringTok{"1"}  \CommentTok{\# E[x**4{-}2*x**2]}
\end{Highlighting}
\end{Shaded}

Matrix path coefficients must be pure rank sums. Mixed fixed-direction
derivatives can be constructed with successive \texttt{directional}
calls and explicitly frozen direction expressions.

\subsection{Gradients, ambient updates and history
pullbacks}\label{gradients-ambient-updates-and-history-pullbacks}

For scalar \texttt{loss},
\texttt{p.gradient(loss,\ vectors={[}x{]},\ matrices={[}W{]},\ scalars={[}s{]})}
returns \texttt{n*grad\_x(loss)}, the ordinary Frobenius gradient
\texttt{grad\_W(loss)}, and ordinary scalar \texttt{grad\_s(loss)}. A
matrix gradient is a finite sum of normalized outer products.
Contributions from every shared forward and transpose use are
accumulated. Additional fixed block multipliers can be included
explicitly in a moving flow, such as
\texttt{\{x:\ -kx*g{[}x{]},\ W:\ -kW*g{[}W{]}\}}.

\texttt{p.gradient\_descent} computes the ambient gradient of the
current-parameter loss template before substituting updated states.
Updates are simultaneous; vector mobility is \texttt{n} and matrix
mobility is one. Its \texttt{mobilities} dictionary supplies extra fixed
deterministic factors. \texttt{steps} is a fixed nonnegative integer;
the step size and factors are deterministic scalar expressions. They are
not allowed to depend on random coordinates or scalar averages. The
helper does not assert loss decrease or step stability.

\begin{Shaded}
\begin{Highlighting}[]
\ImportTok{from}\NormalTok{ pde.mfp\_compiler }\ImportTok{import}\NormalTok{ Program}

\NormalTok{p }\OperatorTok{=}\NormalTok{ Program()}
\NormalTok{lower, upper }\OperatorTok{=}\NormalTok{ p.vector\_type(}\StringTok{"lower"}\NormalTok{), p.vector\_type(}\StringTok{"upper"}\NormalTok{)}
\NormalTok{x }\OperatorTok{=}\NormalTok{ p.root(}\StringTok{"x"}\NormalTok{, lower)}
\NormalTok{W }\OperatorTok{=}\NormalTok{ p.matrix(}\StringTok{"W"}\NormalTok{, lower, upper)}
\NormalTok{loss }\OperatorTok{=}\NormalTok{ p.mean((W }\OperatorTok{@}\NormalTok{ x)}\OperatorTok{**}\DecValTok{2}\NormalTok{)}
\NormalTok{state }\OperatorTok{=}\NormalTok{ p.gradient\_descent(loss, vectors}\OperatorTok{=}\NormalTok{[x], matrices}\OperatorTok{=}\NormalTok{[W], steps}\OperatorTok{=}\DecValTok{2}\NormalTok{, step\_size}\OperatorTok{=}\FloatTok{0.01}\NormalTok{)}
\NormalTok{updated\_loss }\OperatorTok{=}\NormalTok{ p.at(loss, state)}
\BuiltInTok{print}\NormalTok{(p.}\BuiltInTok{compile}\NormalTok{(updated\_loss).formula())}
\end{Highlighting}
\end{Shaded}

\texttt{at} substitutes all supplied states simultaneously without
revisiting their replacement expressions or changing existing frozen
seed expressions. To differentiate through training history, build
\texttt{updated\_loss} first and then differentiate that entire graph.
This pullback includes the parameter dependence of the updates; it is
distinct from the ambient gradient used to take each step. Explicitly
frozen independent seeds remain fixed in that pullback as well.

The \texttt{gradient\_update} calculus example uses the explicit scalar
loss \texttt{mean(z**4)/4}, where \texttt{z=W@one}, a matrix mobility of
one, and fixed step \texttt{eta}. It evaluates the updated transpose on
the saved pre-update vector \texttt{z**3}. The result is exactly lowered
at finite width, and its limiting squared normalized norm is
\texttt{15+(3-15*eta)**2}. A separate \texttt{rank\_update} example
checks \texttt{W+2*u*x.T/n}; \texttt{singular\_queries} uses two calls
\texttt{W@one} with identical inputs. Their source covariance has rank
one, but both formal source slots and both response contributions remain
present.

\subsection{A parameter derivative of a two-hidden-layer
MLP}\label{a-parameter-derivative-of-a-two-hidden-layer-mlp}

The executable \href{code/scripts/example_mfp_mlp_derivative.py}{MLP
example} has input \texttt{x=1}, no biases, two hidden layers of width
\texttt{n}, and stored parameters \texttt{W1,W3} of length \texttt{n},
with \texttt{W2} an \texttt{n} by \texttt{n} matrix. Initialization
entries are independent: \texttt{W1\_i,W3\_i\ \textasciitilde{}\ N(0,1)}
and \texttt{W2\_ij\ \textasciitilde{}\ N(0,1/n)}. This explicitly uses
an order-one stored readout, not the small-readout initialization.

Its forward pass is \texttt{z1=W1}, \texttt{h1=phi(z1)},
\texttt{z2=W2@h1}, \texttt{h2=phi(z2)}, and \texttt{f=W3.T@h2/n}. The
requested first-layer observable is
\texttt{O=n*\textbar{}\textbar{}grad\_W1\ f\textbar{}\textbar{}²},
formed as the normalized squared norm of the returned scaled vector
gradient. The executable constructs the forward pass with
\texttt{Program}, then calls \texttt{gradient}, \texttt{inner}, and
\texttt{compile} as above.

Let \texttt{Z1\ \textasciitilde{}\ N(0,1)},
\texttt{q1=E{[}phi(Z1)**2{]}}, and
\texttt{Z2\ \textasciitilde{}\ N(0,q1)}. Set
\texttt{s1=E{[}phi\textquotesingle{}(Z1)**2{]}},
\texttt{s2=E{[}phi\textquotesingle{}(Z2)**2{]}}, and
\texttt{q2=E{[}phi(Z2)**2{]}}. The limiting expectation of \texttt{O} is
\texttt{s1*s2}. The average first-layer derivative has limit zero; the
full parameter metric norm has limit \texttt{s1*s2+q1*s2+q2}. The reused
transpose is included: its response coefficient vanishes because it
contains the independent centered readout.

{\def\LTcaptype{none} % do not increment counter
\begin{longtable}[]{@{}lrr@{}}
\toprule\noalign{}
Activation flag & First-layer derivative energy & Full metric norm \\
\midrule\noalign{}
\endhead
\bottomrule\noalign{}
\endlastfoot
\texttt{identity} & 1 & 3 \\
\texttt{cubic} & 164025 & 305775 \\
\end{longtable}
}

The \texttt{symbolic} default returns the explicit Gaussian moments. The
test suite checks rational-array backpropagation and the norm identity
independently of the symbolic AD code, then checks both polynomial
specializations.

\subsection{Geometry, labels and feature-kernel
jets}\label{geometry-labels-and-feature-kernel-jets}

The \href{code/scripts/example_mfp_kernel_jets.py}{kernel example} fixes
\texttt{m=d=2} and normalized data
\begin{verbatim}
x1/sqrt(2)=(alpha,0),  x2/sqrt(2)=(beta,gamma).
\end{verbatim}
Geometry and labels \texttt{y1,y2}
remain fixed symbolic scalar parameters. The input Gram is

\begin{Shaded}
\begin{Highlighting}[]
\NormalTok{G11=alpha**2,  G12=alpha*beta,  G22=beta**2+gamma**2.}
\end{Highlighting}
\end{Shaded}

This parameterization is positive semidefinite, including singular
cases. It covers every real PSD two-by-two Gram after choosing factors;
computing factors from arbitrary symbolic covariance entries is not an
API operation.

The shared first-layer matrix has independent columns
\texttt{U,V\ \textasciitilde{}\ N(0,I\_n)}. The hidden matrix
\texttt{W2} has iid \texttt{N(0,1/n)} entries, and the independent
stored readout \texttt{W3} again has iid \texttt{N(0,1)} entries. There
are no biases. The finite forward pass and loss are

\begin{Shaded}
\begin{Highlighting}[]
\NormalTok{z1 = [alpha*U, beta*U+gamma*V]}
\NormalTok{h1[a] = phi(z1[a])}
\NormalTok{z2[a] = W2 @ h1[a]}
\NormalTok{h2[a] = phi(z2[a])}
\NormalTok{f[a] = W3.T @ h2[a] / n}
\NormalTok{r[a] = f[a]{-}y[a]}
\NormalTok{loss = (r[0]**2+r[1]**2)/4.}
\end{Highlighting}
\end{Shaded}

Physical gradient flow trains all four blocks: \texttt{U,V,W3} have
mobility \texttt{n}, and \texttt{W2} has mobility one. The loss is the
\textbf{half mean squared loss}, so its clock differs by a factor of two
from the unhalved loss used in the general neural model. The observable
\texttt{K12=h2{[}0{]}.T@h2{[}1{]}/n} is the last-layer feature kernel,
or equivalently the readout block of the tangent kernel. It is not the
sum over parameter blocks. The script computes

\begin{Shaded}
\begin{Highlighting}[]
\NormalTok{J\_r = lim\_(n→∞) E[(d**r/dt**r) K12(0)] / r!,   r=0,1,2,}
\end{Highlighting}
\end{Shaded}

by taking finite-width moving derivatives first. The residuals, all
parameter blocks and the reused transpose move. A finite-width first
derivative need not vanish. The limiting first coefficient does vanish
in this example.

For generic activation, let \texttt{(Z1\_1,Z1\_2)} be centered Gaussian
with covariance \texttt{G}, define
\texttt{S\_ab=E{[}phi(Z1\_a)*phi(Z1\_b){]}}, and let
\texttt{(Z2\_1,Z2\_2)} have covariance \texttt{S}. Then
\texttt{J0=E{[}phi(Z2\_1)*phi(Z2\_2){]}}, \texttt{J1=0}, and \texttt{J2}
is homogeneous quadratic in \texttt{y1,y2}, with coefficients given by
the generated Gaussian graph. The generic graph is intentionally emitted
by the producer rather than pasted into this guide. Its intermediate
expectations may remain after final output simplification; the graph is
not claimed to be minimal.

Identity activation gives the compact independent check

\begin{Shaded}
\begin{Highlighting}[]
\NormalTok{J0 = G12,  J1 = 0,}
\NormalTok{J2 = (9/4)*(G11*y1+G12*y2)*(G12*y1+G22*y2).}
\end{Highlighting}
\end{Shaded}

For the nonlinear choice \texttt{phi(z)=z**2}, put \texttt{s=G11*G22},
\texttt{c=G12}, and define

\begin{Shaded}
\begin{Highlighting}[]
\NormalTok{A = 5815*s**2+11746*s*c**2+16702*c**4}
\NormalTok{B = 4542*s**4+14636*s**3*c**2+30788*s**2*c**4+13856*s*c**6+4704*c**8.}
\NormalTok{J0 = 9*s**2+2*(s+2*c**2)**2}
\NormalTok{J1 = 0}
\NormalTok{J2 = A*(G11**4*y1**2+G22**4*y2**2)+B*y1*y2.}
\end{Highlighting}
\end{Shaded}

The complete producer regenerates that polynomial with symbolic geometry
and labels. Tests compare its compact form to the generated expansion,
and also specialize every generic expectation atom to the quadratic
activation and compare with direct polynomial compilation. Those two
compilation paths share the source-rule implementation; their agreement
is not an independent proof of the full second coefficient. The
finite-width AD check is independent: rational dense backpropagation and
coefficient convolution through order two verify all moving blocks at
widths two and three. The initial nonlinear kernel is separately checked
by two Gaussian fourth-moment identities. These checks require no
retained outputs.

\subsection{Returned graph, expression API and finite
interpreter}\label{returned-graph-expression-api-and-finite-interpreter}

\texttt{dag.expectations} is an ordered tuple of entries with fields
\texttt{symbol}, \texttt{integrand}, \texttt{coordinates},
\texttt{covariance}, and \texttt{kind}. Covariance entries and
integrands use only prior expectation values and fixed deterministic
parameters. \texttt{dag.output} is the final scalar expression.
\texttt{dag.fields} exposes computed representatives;
\texttt{dag.sources} retains named source coordinates and covariances,
even for duplicate calls with perfect correlation. \texttt{dag.trace}
includes the finite builder transformations and Gaussian dependency
traversal. \texttt{formula()} renders definitions once;
\texttt{to\_dict()} yields a JSON-ready description. These are
inspection outputs, not a serializer for recreating an editable
\texttt{Program}.

Two distinct derivative engines are exposed.
\texttt{Program.directional} and \texttt{Program.gradient} propagate
through the finite typed graph, including scalar feedback and reused
parameters. \texttt{pde.mfp\_expr.diff} takes named scalar-source
partials inside Gaussian evaluation, holding all other symbols,
including previous expectation coefficients, fixed. Singular covariance
does not identify formal variables before these partials are taken.

The supporting \texttt{pde.mfp\_expr} factories are \texttt{const},
\texttt{symbol}, and \texttt{phi}; its operations include \texttt{diff},
\texttt{expand}, \texttt{substitute}, \texttt{evaluate},
\texttt{render}, and \texttt{gaussian\_expectation}. Substitution is
simultaneous and structural.
\texttt{evaluate(expr,\ values,\ activation)} calls
\texttt{activation(order,\ value)} for activation nodes; it evaluates a
supplied scalar expression, not an entire expectation DAG.
\texttt{gaussian\_expectation} accepts named centered coordinates, with
missing covariance entries treated as zero. Symbolic covariance symmetry
is checked structurally, while positive semidefiniteness is the caller's
obligation at this lower level. Polynomial terms reduce by Wick's rule;
flat activation products use terminating Stein reduction. Nested
activation integrals remain atoms. Rational moment checks compare the
specialized reducer with the maintained
\texttt{pde.gaussian\_moments.gaussian\_moment} API, whose stricter
no-float input contract is unchanged.

The independent interpreter is
\texttt{pde.mfp\_finite.evaluate\_finite(output,\ width,\ roots,\ matrices,\ parameters=None,\ activation=None)}.
It evaluates actual supplied arrays and never samples the initialization
law. Pass finite Python lists/tuples and actual stored matrix entries,
with no extra matrix rescaling. All required root/matrix/parameter
leaves need values; unused supplied entries are validated too. The
activation callback must supply the actual derivative at its requested
order. Integer and \texttt{Fraction} inputs remain exact; float state
entries use ordinary floating arithmetic, unlike symbolic
decimal-rational conversion. Vector outputs are immutable tuples. Finite
input checks do not guarantee that every intermediate floating operation
avoids overflow or underflow.

\begin{Shaded}
\begin{Highlighting}[]
\ImportTok{from}\NormalTok{ fractions }\ImportTok{import}\NormalTok{ Fraction}
\ImportTok{from}\NormalTok{ pde.mfp\_compiler }\ImportTok{import}\NormalTok{ Program}
\ImportTok{from}\NormalTok{ pde.mfp\_finite }\ImportTok{import}\NormalTok{ evaluate\_finite}

\NormalTok{p }\OperatorTok{=}\NormalTok{ Program()}
\NormalTok{x }\OperatorTok{=}\NormalTok{ p.root(}\StringTok{"x"}\NormalTok{, p.vector\_type(}\StringTok{"neurons"}\NormalTok{))}
\NormalTok{O }\OperatorTok{=}\NormalTok{ p.mean(x}\OperatorTok{**}\DecValTok{2}\NormalTok{)}
\ControlFlowTok{assert}\NormalTok{ evaluate\_finite(O, }\DecValTok{2}\NormalTok{, \{x: [Fraction(}\DecValTok{1}\NormalTok{, }\DecValTok{2}\NormalTok{), Fraction(}\DecValTok{3}\NormalTok{, }\DecValTok{2}\NormalTok{)]\}, \{\}) }\OperatorTok{==}\NormalTok{ Fraction(}\DecValTok{5}\NormalTok{, }\DecValTok{4}\NormalTok{)}
\end{Highlighting}
\end{Shaded}

No part of this API silently replaces a Gaussian integral with numerical
quadrature, supplies a covariance family with a neural interpretation,
or turns fixed finite derivatives into a positive-time population
solver.

\section{Source bundle and reproduction}\label{sec-source-bundle}

The standalone directory supplies \texttt{main.tex},
\texttt{references.bib}, a \texttt{Makefile}, and the accompanying
\texttt{code/} tree. The native LaTeX source compiles with XeLaTeX,
BibTeX and \texttt{latexmk}; Quarto is not needed to build this paper.
From the paper directory, run \texttt{make\ pdf}. The output is
\texttt{main.pdf}. Install the Python dependency with
\texttt{python\ -m\ pip\ install\ -r\ requirements.txt} in an
environment of your choice, then run \texttt{make\ check} for the copied
test suite.

The symbolic implementation is in \texttt{code/pde/mfp\_compiler.py} and
\texttt{code/pde/mfp\_expr.py}; the independent finite-array interpreter
is in \texttt{code/pde/mfp\_finite.py}. The package initialization,
finite-network module and exact Gaussian-moment module are included to
preserve the established import closure. Three
\texttt{code/scripts/example\_mfp\_*.py} files implement the worked
calculations, and four \texttt{code/tests/test\_mfp\_*.py} files supply
the established checks. All Python sources and the dependency pin are
copied byte for byte from the maintained package.
\texttt{SOURCE\_MAP.json} records the source hashes, manuscript
extraction ranges and editorial adaptations.

\appendix

\section{Bounded-derivative Gaussian-program
foundation}\label{sec-docs-special-data-limits-l3785}

This part treats every fixed finite hidden depth L. All instruction
lists and the number of initialized matrices are fixed before width
tends to infinity. Its elementary proofs do not assert uniformity for a
depth or transcript length growing with width.

\subsection{Finite programs and convergence of their empirical
laws}\label{sec-docs-special-data-limits-l3789}

There are \(L\) types of length-\(n\) vectors, one for each hidden
layer. Operations combining coordinates may combine only vectors of the
same type. For each \(2\le\ell\le L\), let

\[
 W^{(\ell)}_n:\mathbb R^n_{\ell-1}\longrightarrow\mathbb R^n_\ell
\]

be mutually independent matrices with independent \(N(0,1/n)\) entries.
Their transposes are reused as the reverse actions of these same
matrices.

Each layer may have a fixed finite tuple of root vectors. Its coordinate
tuples are independent and identically distributed, have finite second
moment, and are independent of all matrices. Tuples in different layers
are independent. Constants are also allowed. Additional independent
Gaussian roots may be added to any layer when a proof requires probes or
regularization.

A deterministic-coefficient program is a fixed finite ordered list of
instructions of the following forms:

\begin{enumerate}
\def\labelenumi{\arabic{enumi}.}
\tightlist
\item
  apply a fixed \(C^1\) function \(F:\mathbb R^m\to\mathbb R\) with
  bounded first partial derivatives, coordinate by coordinate, to
  previously available vectors of one layer;
\item
  multiply a previously available vector by \(W^{(\ell)}_n\) or
  \({W^{(\ell)}_n}^{\top}\) for any \(2\le\ell\le L\), with the
  appropriate types;
\item
  form a fixed real linear combination of previous same-layer vectors.
\end{enumerate}

The first condition implies a global Lipschitz bound and at most linear
growth for each coordinate instruction. The bound may depend on that
instruction. Root tuples themselves need not be generated by such
functions.

Write

\[
\frac{1}{n}\langle u,v\rangle_{\mathbb R^n}=\frac1n\sum_{\alpha=1}^n u_\alpha v_\alpha,
\qquad \frac{\|u\|_2^2}{n}=\frac{1}{n}\langle u,u\rangle_{\mathbb R^n}.
\]

For same-layer nodes \(v^1_n,\ldots,v^m_n\), their empirical law is

\[
\widehat\mu_n=\frac1n\sum_{\alpha=1}^n
 \delta_{(v^1_{n,\alpha},\ldots,v^m_{n,\alpha})}.
\]

Here \(\mathcal W_2\) uses the Euclidean distance on \(\mathbb R^m\).

\begin{theorem}[fixed finite Gaussian
program]\protect\hypertarget{thm-docs-special-data-limits-l3820}{}\label{thm-docs-special-data-limits-l3820}

Every such program has deterministic joint limiting laws of all its
same-layer node tuples, and

\begin{equation}\protect\phantomsection\label{eq-docs-special-data-limits-l3821}{
\mathcal W_2(\widehat\mu_n,\mu)\longrightarrow0
\quad\hbox{in probability}
}\end{equation}

along the full width sequence. In particular every within-layer pairwise
contraction converges to the corresponding limiting second moment. The
scalar laws are given by the source rule in
Section~\ref{sec-docs-special-data-limits-l3946}. Query Grams may be
singular. Finite collections of programs sharing the same matrices and
roots converge jointly by applying the assertion to their finite union.

We prove this theorem in Sections
\ref{sec-docs-special-data-limits-l3842}--\ref{sec-docs-special-data-limits-l4009}.
The following elementary facts make explicit the probabilistic mode of
convergence used in its proof.

If \(X_\alpha\) are iid and \(E|X_1|<\infty\), their averages converge
in probability to their expectation: truncate \(X_\alpha\) at level
\(M\), use the variance bound \(O(M^2/n)\) for the bounded variables,
and bound the mean absolute truncation error by \(E[|X_1|1_{|X_1|>M}]\).
First send \(n\) to infinity and then \(M\) to infinity. This proves the
required initial weak convergence and second-moment convergence of root
empirical laws.

For probability measures on a finite-dimensional Euclidean space, weak
convergence together with convergence of second moments implies
\(\mathcal W_2\) convergence. One direct proof is as follows. Continuous
truncations of \(|x|^2\) show that the second moments outside
sufficiently large balls are uniformly small. Inside a ball partition
space into finitely many sets of diameter at most \(\eta\), choosing
boundaries of zero limiting measure. Weak convergence makes their masses
converge. Couple the common mass within each partition cell, at cost at
most \(\eta^2\), and couple the remaining masses arbitrarily. The
unmatched mass inside the ball vanishes; its cost is bounded by the
squared diameter of the ball times that mass. The tails have arbitrarily
small cost by \(|x-y|^2\le2|x|^2+2|y|^2\). Sending the ball radius and
then the partition resolution to their limits proves the claim. A
countable family of bounded Lipschitz tests determines weak convergence,
by approximation on compact balls and tightness. For random measures the
same argument applies in probability, or along an almost surely
convergent subsubsequence of every subsequence.

Two arrays on the same neuron indices satisfy

\begin{equation}\protect\phantomsection\label{eq-docs-special-data-limits-l3835}{
\mathcal W_2^2(\widehat\mu_n,\widehat\nu_n)
\le\frac1n\sum_\alpha |X_{n,\alpha}-Y_{n,\alpha}|^2,
}\end{equation}

using the coupling that pairs equal indices. These facts require no
assertion that trained coordinates are independent.

\end{theorem}

\subsection{An explicit Gaussian operator norm
bound}\label{sec-docs-special-data-limits-l3842}

\begin{lemma}[]\protect\hypertarget{lem-docs-special-data-limits-l3844}{}\label{lem-docs-special-data-limits-l3844}

For an \(n\times n\) matrix \(W_n\) with independent \(N(0,1/n)\)
entries,

\begin{equation}\protect\phantomsection\label{eq-docs-special-data-limits-l3845}{
\Pr(\|W_n\|_{\mathrm{op}}>10)
\le 2\,9^{2n}e^{-100n/8}\longrightarrow0.
}\end{equation}

\end{lemma}

\protect\phantomsection\label{proof-docs-special-data-limits-l3851}{}

\begin{proof}
A maximal \(1/4\)-separated subset \(\mathcal N\) of the Euclidean unit
sphere is a \(1/4\)-net. Balls of radius \(1/8\) about its points are
disjoint and lie in the ball of radius \(9/8\), so volume comparison
gives \(|\mathcal N|\le9^n\). For unit \(u,v\), choose
\(u_0,v_0\in\mathcal N\) within \(1/4\). Then

\[
|u^TW_nv-u_0^TW_nv_0|
\le\tfrac12\|W_n\|_{\mathrm{op}}.
\]

Taking the supremum gives
\(\|W_n\|_{\mathrm{op}}\le2\max_{u_0,v_0\in\mathcal N}|u_0^TW_nv_0|\).
For a fixed pair the displayed scalar is \(N(0,1/n)\). Its exponential
moment is \(Ee^{t u_0^TW_nv_0}=e^{t^2/(2n)}\); Markov's inequality
optimized at \(t=ns\) yields \(\Pr(|u_0^TW_nv_0|>s)\le2e^{-ns^2/2}\). A
union bound at \(s=5\) proves (\ref{eq-docs-special-data-limits-l3845}).
The exponent is negative since \(2\log9<12.5\). The same bound applies
to transposes, and a finite union bound handles all matrices.
\(\blacksquare\)
\end{proof}

The same argument for a threshold \(t\ge10\) gives

\begin{equation}\protect\phantomsection\label{eq-docs-special-data-limits-l3859}{
\Pr(\|W_n\|_{\mathrm{op}}>t)
\le2\exp\{n(2\log9-t^2/8)\}
\le2e^{-nt^2/16}\le2e^{-t^2/16}.
}\end{equation}

Consequently every fixed positive moment is bounded uniformly in width:

\begin{equation}\protect\phantomsection\label{eq-docs-special-data-limits-l3866}{
\sup_{n\ge1}E\|W_n\|_{\mathrm{op}}^p
\le10^p+2p\int_{10}^\infty t^{p-1}e^{-t^2/16}\,dt<\infty.
}\end{equation}

The integration formula follows by writing \(X^p=\int_0^Xpt^{p-1}dt\)
for nonnegative \(X\) and interchanging nonnegative integrals. Hölder's
inequality then gives uniform fixed-order moments for every fixed
polynomial in finitely many such operator norms. For a standard Gaussian
vector \(g_n\), Jensen's inequality also gives
\(E\frac{\|g_n\|_2^p}{n^{p/2}}\le E|G|^p\) when \(p\ge2\).

A useful consequence identifies normalized traces without a
concentration theorem for functions of matrix entries. Let \(T_n\) be
any random real \(n\times n\) matrix independent of
\(g_n\sim N(0,I_n)\), with \(\sup_nE\|T_n\|_{\mathrm{op}}^2<\infty\).
Conditional on \(T_n\),

\begin{equation}\protect\phantomsection\label{eq-docs-special-data-limits-l3874}{
E_g\frac{1}{n}\langle g_n,T_ng_n\rangle_{\mathbb R^n}=\frac1n\operatorname{tr}T_n,
\quad
\operatorname{Var}_g\!\left(\frac{1}{n}\langle g_n,T_ng_n\rangle_{\mathbb R^n}\right)
=\frac{2}{n^2}\left\|\frac{T_n+T_n^T}{2}\right\|_F^2
\le\frac{2}{n}\|T_n\|_{\mathrm{op}}^2.
}\end{equation}

To verify the variance, replace \(T_n\) by its symmetric part,
diagonalize it orthogonally, and use that the transformed Gaussian
coordinates are independent with \(\operatorname{Var}(G^2)=2\).
Therefore the difference between this probe and the normalized trace
tends to zero in \(L^2\). If the probe is a fixed finite program,
Theorem~\ref{thm-docs-special-data-limits-l3820} identifies its
deterministic limit and hence the trace limit in probability. A fixed
polynomial in the initialized actions and their adjoints satisfies the
operator moment hypothesis by (\ref{eq-docs-special-data-limits-l3866}),
and its application to the probe is such a program. Uniform moments of
order greater than one upgrade convergence of these normalized traces to
convergence of their expectations: split at a large absolute threshold
and use the higher-moment bound to make the first-moment tails uniformly
small. The same observation supplies uniform integrability of every
fixed polynomial expression needed for finite-degree moment
calculations.

\subsection{Exact adaptive Gaussian
conditioning}\label{sec-docs-special-data-limits-l3884}

Let the current transcript consist of all revealed roots and all
previously computed vectors. For one matrix \(W\), collect its earlier
forward and reverse observations as

\begin{equation}\protect\phantomsection\label{eq-docs-special-data-limits-l3887}{
WV=Y,\qquad W^TU=Q.
}\end{equation}

The columns of \(V\) and \(U\) are the respective query inputs, with
output columns in \(Y\) and \(Q\). Conditioned on the transcript they
are fixed. Empty column lists are allowed; terms involving them are
omitted.

It is necessary to justify this conditioning for adaptive inputs.
Initially the conditional laws of the matrices are independent Gaussian
laws. Suppose this is true, with the linear constraints already
observed, at a particular instruction. A coordinate operation is
measurable from the transcript and reveals no new randomness. At a
matrix call its input is also measurable from the transcript.
Conditional on the transcript, the new answer is a linear observation of
only the queried matrix. Conditioning a product of the current
conditional laws on this observation leaves the other factors unchanged
and conditions only the queried factor. Thus induction preserves
independence of the residual matrix factors and adds exactly the
indicated linear constraint. A freshly revealed independent root
likewise does not alter these residual laws. This argument conditions
successively, and does not assume that an adaptive input was independent
of the matrix before the transcript was fixed.

Suppose first that \(V^TV\) and \(U^TU\) are invertible. Let
\(P_V=V(V^TV)^{-1}V^T\), and similarly define \(P_U\). Then

\begin{equation}\protect\phantomsection\label{eq-docs-special-data-limits-l3896}{
W\mid\mathcal H\ \overset d=
M+P_{U^\perp}\widetilde W P_{V^\perp},
\quad
M=Y(V^TV)^{-1}V^T
 +U(U^TU)^{-1}Q^TP_{V^\perp},
}\end{equation}

where \(\widetilde W\) is an independent copy of the original matrix.

Here is a direct verification of the Gaussian projection behind
(\ref{eq-docs-special-data-limits-l3896}). The compatibility relation is
\(U^TY=Q^TV\), because both sides equal \(U^TWV\). It gives \(MV=Y\) and
\(M^TU=Q\). The homogeneous solutions of
(\ref{eq-docs-special-data-limits-l3887}) are exactly matrices
\(K=P_{U^\perp}KP_{V^\perp}\). Both summands defining \(M\) are
orthogonal in Frobenius inner product to that subspace. Hence \(M\) is
the unique minimum-Frobenius-norm solution. Vectorize \(W\), whose law
is an isotropic Gaussian in \(\mathbb R^{n^2}\). In an orthonormal basis
adapted to the homogeneous solution subspace its coordinates are
independent Gaussians; conditioning on the orthogonal coordinates leaves
independent Gaussians on the homogeneous subspace and fixes the other
coordinates to those of \(M\). This proves
(\ref{eq-docs-special-data-limits-l3896}). It also proves the same
assertion with orthogonal projections and minimum-norm solutions when
column lists are linearly dependent; the nonsingular formula is the only
one whose coefficients we take to a width limit.

For a new forward input \(h\), put

\[
\alpha_n=(V^TV)^{-1}V^Th,
\quad h_\perp=h-V\alpha_n,
\quad
\beta_n=(U^TU/n)^{-1}(Q^Th_\perp/n).
\]

Equation (\ref{eq-docs-special-data-limits-l3896}) becomes

\begin{equation}\protect\phantomsection\label{eq-docs-special-data-limits-l3916}{
Wh=Y\alpha_n+U\beta_n
 +\frac{\|h_\perp\|_2}{\sqrt n}P_{U^\perp}g
\quad\hbox{in conditional law},
}\end{equation}

with \(g\sim N(0,I_n)\) independent of the transcript. The reverse
formula follows by interchanging the two sides.

Assume provisionally that every query Gram has a positive definite
limit. Inductively all contractions of old nodes converge. Thus the
coefficients in (\ref{eq-docs-special-data-limits-l3916}), and
\(\frac{\|h_\perp\|_2}{\sqrt n}\), converge in probability to
deterministic limits. Inverting a positive definite fixed-size matrix is
continuous, for instance by a Neumann-series expansion about its
invertible limit.

The projection removed from the fresh noise is negligible:

\begin{equation}\protect\phantomsection\label{eq-docs-special-data-limits-l3927}{
E[\frac{\|P_Ug\|_2^2}{n}\mid\mathcal H]
=\frac{\operatorname{rank}U}{n}.
}\end{equation}

The multiplying variance factor is bounded in probability, so
conditional Markov's inequality makes its contribution vanish in
normalized mean square. After this removal and replacement of convergent
coefficients by their limits, the new coordinate is a deterministic
linear combination \(m_\alpha\) of old same-layer nodes plus
\(\sigma g_\alpha\).

For a bounded Lipschitz test \(\psi\) of the old tuple and this new
coordinate, conditional independence of \(g_\alpha\) gives variance at
most \(4\|\psi\|_\infty^2/n\) for its empirical average. Its conditional
mean is the old empirical average of the bounded continuous function

\[
x\longmapsto E_G\psi(x,m(x)+\sigma G),
\]

which converges by the induction hypothesis. For the new second moment
expand

\[
\frac1n\sum_\alpha(m_\alpha+\sigma g_\alpha)^2
=\frac{\|m\|_2^2}{n}+\frac{2\sigma}{n}\sum_\alpha m_\alpha g_\alpha
 +\frac{\sigma^2}{n}\sum_\alpha g_\alpha^2.
\]

The middle term has conditional variance
\(4\sigma^2\frac{\|m\|_2^2}{n}/n\), and the last average has variance
\(2/n\). All relevant norms are bounded in probability. We obtain weak
convergence and second-moment convergence, hence
(\ref{eq-docs-special-data-limits-l3821}). Coordinate instructions
preserve this convergence because their Lipschitz constants bound the
transport cost. This completes the induction under the provisional
positive-definiteness assumption.

\subsection{Source-response identity and formal
derivatives}\label{sec-docs-special-data-limits-l3946}

For every oriented initialized matrix introduce a centered Gaussian
source group indexed by its calls. Sources for different orientations,
including a matrix and its transpose, are independent groups; they are
also independent of the root tuples. Within a forward group for \(W\),
the source attached to input \(h\) has covariance with the source
attached to input \(v\) equal to \(E[hv]\). Within the reverse group the
analogous covariance is \(E[uv]\) for the corresponding reverse inputs.
These are uncentered second moments of inputs and centered covariances
of sources.

The scalar node of a new forward call is

\begin{equation}\protect\phantomsection\label{eq-docs-special-data-limits-l3951}{
\mathscr W h=\xi_h+\sum_{s:\,W^Tu_s\text{ already called}}
 u_s\,E[\partial_{\zeta_s}h].
}\end{equation}

The scalar node of a reverse call is

\begin{equation}\protect\phantomsection\label{eq-docs-special-data-limits-l3957}{
\mathscr W^*u=\zeta_u+\sum_{r:\,Wv_r\text{ already called}}
 v_r\,E[\partial_{\xi_r}u].
}\end{equation}

An input is its explicit expression in named source coordinates and
roots, obtained by unrolling previous scalar instructions. A derivative
in (\ref{eq-docs-special-data-limits-l3951}) or
(\ref{eq-docs-special-data-limits-l3957}) differentiates that
expression. Previously computed expectations, coefficients, covariance
entries, mesh sizes, and any deterministic control values are held
fixed. Each named source remains a separate formal argument, including
when the joint source covariance is singular. An unavailable source has
derivative zero. Derivative paths through other matrices' earlier calls
remain part of the expression.

The recursion is causal. At a call, its input and its source derivatives
are already defined; their expectations determine the response
coefficients. The source covariance extension is the Gram extension of
the corresponding input list and is therefore positive semidefinite. A
Gaussian group with that extended covariance exists: if the old
covariance is \(K\), the new cross-covariance is \(b\), and the new
variance is \(v\), positivity implies \(b\in\operatorname{ran}K\) and
\(v-b^TK^+b\ge0\). To see the range assertion, test positivity on
\((tu,1)\) with \(Ku=0\) and arbitrary \(t\). Completing the square on
\(\operatorname{ran}K\) gives the second assertion. Consequently the new
coordinate can be represented as \(b^TK^+\xi+\sqrt{v-b^TK^+b}\,G\), with
a fresh standard normal \(G\). Here a pseudoinverse is used only to
construct one fixed finite Gaussian law; no continuity of pseudoinverses
is asserted.

\begin{lemma}[source
rule]\protect\hypertarget{lem-docs-special-data-limits-l3966}{}\label{lem-docs-special-data-limits-l3966}

Under the positive-definiteness assumption of
Section~\ref{sec-docs-special-data-limits-l3884},
(\ref{eq-docs-special-data-limits-l3951}) and
(\ref{eq-docs-special-data-limits-l3957}) give exactly the scalar laws
obtained there.

\end{lemma}

\protect\phantomsection\label{proof-docs-special-data-limits-l3968}{}

\begin{proof}
Consider a forward call and use the notation of
(\ref{eq-docs-special-data-limits-l3916}). Write old forward inputs as
\(v_r\), old reverse inputs as \(u_s\), and their scalar outputs, by
induction, as

\[
y_r=\xi_r+\sum_s D_{rs}u_s,
\qquad D_{rs}=E[\partial_{\zeta_s}v_r],
\]

\[
q_s=\zeta_s+\text{a deterministic linear combination of old }v_r.
\]

All old source lists are padded by zeros for unavailable indices. Let
\(\alpha\) be the limiting least-squares coefficient from
(\ref{eq-docs-special-data-limits-l3916}), and let
\(h_\perp=h-\sum_r\alpha_rv_r\). Orthogonality gives
\(E[v_rh_\perp]=0\). Therefore

\begin{equation}\protect\phantomsection\label{eq-docs-special-data-limits-l3977}{
E[q_sh_\perp]=E[\zeta_sh_\perp].
}\end{equation}

Let \(G_U=(E[u_su_t])_{st}\), the covariance matrix of \(\zeta\).
Gaussian integration by parts gives

\begin{equation}\protect\phantomsection\label{eq-docs-special-data-limits-l3982}{
E[\zeta h_\perp]=G_U E[\nabla_\zeta h_\perp].
}\end{equation}

For completeness, the one-dimensional identity \(E[Gf(G)]=E[f'(G)]\)
follows by integration by parts against the standard normal density. The
boundary term vanishes for a function of at most linear growth with
bounded derivative. Represent a possibly singular Gaussian vector as
\(\zeta=T G\), apply this identity in each independent standard normal
coordinate of \(G\), and sum using \(TT^T=G_U\). Conditioning on
independent roots and the other source groups proves
(\ref{eq-docs-special-data-limits-l3982}) in the present setting. Every
derivative is integrable: at a fixed finite instruction, its norm is
bounded by a deterministic finite expression in earlier coefficients and
the bounded derivatives of coordinate maps.

The limiting coefficient of \(U\) in
(\ref{eq-docs-special-data-limits-l3916}) is consequently

\[
\beta=E[\nabla_\zeta h]-\sum_r\alpha_r E[\nabla_\zeta v_r].
\]

Substitution of the old \(y_r\) decompositions in
(\ref{eq-docs-special-data-limits-l3916}) cancels the second term
exactly. The answer becomes

\[
\xi_h+\sum_su_sE[\partial_{\zeta_s}h],
\qquad
\xi_h=\sum_r\alpha_r\xi_r+\sigma G,
\quad \sigma^2=E[h_\perp^2].
\]

Since the old \(\xi\) covariance is the Gram of the \(v_r\),

\[
E[\xi_h\xi_r]=E[hv_r],
\quad
E[\xi_h^2]=E\Big(\sum_r\alpha_rv_r\Big)^2+E[h_\perp^2]=E[h^2].
\]

The fresh normal is independent of all old roots and source groups. Thus
adjoining it preserves independence of distinct oriented source groups.
The reverse calculation is the same after interchanging the two layers.
Interleaving calls of different matrices does not alter this
calculation, because Section~\ref{sec-docs-special-data-limits-l3884}
established the conditional independence of their residual factors.
\(\blacksquare\)
\end{proof}

Independence of source groups does not assert that the answers of a
matrix and its transpose are independent. Their response terms encode
their dependence. Nor does it assert that a source is independent of all
later inputs; later scalar inputs can be functions of that source.

\subsection{Singular queries without a rank-stability
assumption}\label{sec-docs-special-data-limits-l4009}

\begin{lemma}[regularization of a fixed
program]\protect\hypertarget{lem-docs-special-data-limits-l4011}{}\label{lem-docs-special-data-limits-l4011}

The conclusions of Theorem~\ref{thm-docs-special-data-limits-l3820} and
the formulas
(\ref{eq-docs-special-data-limits-l3951})--(\ref{eq-docs-special-data-limits-l3957})
hold when any of the limiting input Grams is singular.

\end{lemma}

\protect\phantomsection\label{proof-docs-special-data-limits-l4013}{}

\begin{proof}
For every matrix call introduce a new independent standard Gaussian
input vector \(\chi\), revealed immediately before that call, and
replace its input \(h\) by \(h+\varepsilon\chi\). Each call has a
distinct noise vector. The other instructions are unchanged.

At fixed \(\varepsilon>0\), the new noise is independent of the old
transcript and of the unperturbed part of the current input. If \(V\) is
the list of prior same-orientation inputs, the normalized squared
distance from \(h+\varepsilon\chi\) to \(\operatorname{span}V\) is

\[
\frac{\|P_{V^\perp}h\|_2^2}{n}
 +2\varepsilon\frac{1}{n}\langle P_{V^\perp}h,\chi\rangle_{\mathbb R^n}
 +\varepsilon^2\frac{\|P_{V^\perp}\chi\|_2^2}{n}.
\]

Conditionally, the cross term has variance
\(4\varepsilon^2\frac{\|P_{V^\perp}h\|_2^2}{n}/n\). The last norm
squared has mean \(1-\operatorname{rank}V/n\) and variance at most
\(2/n\). Thus every limiting new squared distance is at least
\(\varepsilon^2\). Induction gives positive definite limiting query
Grams, so Sections
\ref{sec-docs-special-data-limits-l3884}--\ref{sec-docs-special-data-limits-l3946}
apply to the perturbed program. Equivalently, in its scalar law the new
independent root adds \(\varepsilon^2\) to the Schur complement of the
old input Gram.

Couple the perturbed and original finite programs with the same matrices
and roots. On the event that all initialized matrix norms are at most 10
and that all of the finitely many fresh noise vectors have normalized
norms at most 2, propagate errors instruction by instruction. A
coordinate instruction multiplies the previous error by its fixed
Lipschitz constant; a linear combination contributes the sum of
coefficient magnitudes times previous errors; a matrix call contributes
at most ten times the input error plus \(20\varepsilon\). Consequently

\begin{equation}\protect\phantomsection\label{eq-docs-special-data-limits-l4024}{
\max_{\text{nodes }v}\frac{\|v_n^\varepsilon-v_n\|_2}{\sqrt n}
\le C\varepsilon,
}\end{equation}

where \(C\) is finite and independent of \(n\) and
\(0<\varepsilon\le1\). The event has probability tending to one by
Lemma~\ref{lem-docs-special-data-limits-l3844} and the elementary
second-moment calculation for Gaussian noise norms.

We next show that the scalar recursion itself is continuous at
\(\varepsilon=0\); this step concerns covariances and derivatives, not
inverses of empirical Grams. Induct on its finitely many instructions.
Each scalar node is a \(C^1\) expression in the finite named source list
and roots. If earlier deterministic coefficients remain in a compact
set, the expression and its first source derivatives have uniform
bounds: the expression has at most linear growth in the root and source
coordinates, and its derivatives have a finite deterministic bound. This
follows directly by applying the coordinate derivative bounds and the
linear response formulas in the previous instructions. The values and
first derivatives are continuous in their arguments and in the earlier
coefficient list.

By induction, the covariance entries for the next source, which are
second moments of old scalar inputs, converge as
\(\varepsilon\downarrow0\). If positive semidefinite matrices
\(K_j\to K\) have fixed size, then \(K_j^{1/2}\to K^{1/2}\). To verify
this without a regularity assumption on eigenvalues, their positive
square roots are bounded. Every convergent subsequence of these square
roots has a positive semidefinite limit \(T\) with \(T^2=K\). A positive
semidefinite matrix has a unique positive semidefinite square root:
diagonalize it, observe that any such \(T\) commutes with \(K=T^2\), and
restrict to its eigenspaces. Hence every subsequential limit is
\(K^{1/2}\), which proves convergence.

Represent the full finite source prefix for each \(\varepsilon\) as
\(K_\varepsilon^{1/2}G\) using one standard Gaussian vector for each
oriented group, independently of the roots. This couples the source
prefixes in \(L^2\). The uniform linear-growth bounds and Lipschitz
constants for node expressions then give their \(L^2\) convergence. More
explicitly, split the node difference into a change of arguments at
fixed coefficients, bounded by the common Lipschitz constant, and a
change of coefficients at fixed arguments. The latter converges
pointwise and is bounded by a constant times one plus the norm of the
finite root/source list, an \(L^2\) dominator. First source derivatives
converge in probability and are uniformly bounded, so their expectations
converge. This proves convergence of the next response coefficient and
closes the induction. At zero noise the resulting expression is exactly
(\ref{eq-docs-special-data-limits-l3951})--(\ref{eq-docs-special-data-limits-l3957})
for the original formal program.

Let \(\mu^\varepsilon\) be the perturbed scalar law of a selected tuple
and \(\mu^0\) the zero-noise law just constructed. We have
\(\mathcal W_2(\mu^\varepsilon,\mu^0)\to0\). By
(\ref{eq-docs-special-data-limits-l3835}),
(\ref{eq-docs-special-data-limits-l4024}), and the proved
fixed-\(\varepsilon\) limit,

\[
\mathcal W_2(\widehat\mu_n,\mu^0)
\le C_m\varepsilon
 +\mathcal W_2(\widehat\mu_n^\varepsilon,\mu^\varepsilon)
 +\mathcal W_2(\mu^\varepsilon,\mu^0)
\]

on an event of probability tending to one. Choose \(\varepsilon\) first,
let \(n\to\infty\), and then let \(\varepsilon\downarrow0\). This proves
the full-sequence convergence in probability, including singular Grams,
and finishes Theorem~\ref{thm-docs-special-data-limits-l3820}.
\(\blacksquare\)
\end{proof}

The derivative convention has a precise invariant meaning on singular
supports. If a source vector \(\zeta\) has covariance \(G\), and \(u\)
is the vector of the associated reverse inputs with \(E[uu^T]=G\), then

\begin{equation}\protect\phantomsection\label{eq-docs-special-data-limits-l4047}{
E[\zeta f]=G E[\nabla f],\qquad
u^Tv=0\text{ a.s. for every }v\in\ker G.
}\end{equation}

The second identity follows from \(E[(u^Tv)^2]=v^TGv=0\). If two
admissible smooth formal expressions agree on the Gaussian support,
their expected derivative vectors differ by an element of \(\ker G\), by
the first identity. Their contracted corrections therefore agree.
Individual derivative coefficients need not agree. None of this implies
that pseudoinverses converge at rank loss.

\subsection{Causal scalar
feedback}\label{sec-docs-special-data-limits-l4054}

The deterministic-coefficient theorem also identifies programs with the
following causal scalar feedback. At finitely many stages, compute inner
products of already available same-layer nodes, apply locally Lipschitz
functions to the resulting finite scalar list, and use the resulting
numbers as coefficients of subsequent linear combinations. Assume all
scalar operations are defined on a neighborhood of their deterministic
limiting arguments; divisions require a nonzero limiting denominator.
Require the actual finite operation to be defined everywhere it is used,
or assign an arbitrary measurable fallback outside that neighborhood.
Convergence of its arguments makes the exceptional event have
probability tending to zero. The physical algorithms themselves use no
division. Coefficients may multiply unbounded vector nodes, because
their perturbations can be estimated by the vector's normalized \(L^2\)
norm.

To prove this extension, construct an oracle program by replacing each
scalar feedback value by its limiting deterministic value, computed from
earlier scalar nodes. The construction is causal and therefore not an
implicit fixed-point definition.
Theorem~\ref{thm-docs-special-data-limits-l3820} identifies this oracle.
At the next scalar step use

\begin{equation}\protect\phantomsection\label{eq-docs-special-data-limits-l4059}{
|\frac{1}{n}\langle u,v\rangle_{\mathbb R^n}-\frac{1}{n}\langle\bar u,\bar v\rangle_{\mathbb R^n}|
\le\frac{\|u-\bar u\|_2}{\sqrt n}\frac{\|v\|_2}{\sqrt n}
 +\frac{\|\bar u\|_2}{\sqrt n}\frac{\|v-\bar v\|_2}{\sqrt n}.
}\end{equation}

At the next scalar multiplication use

\[
\frac{\|c u-\bar c\bar u\|_2}{\sqrt n}
\le |c|\frac{\|u-\bar u\|_2}{\sqrt n}+|c-\bar c|\frac{\|\bar u\|_2}{\sqrt n}.
\]

All oracle norms and finitely many oracle coefficients are bounded in
probability; initial operator norms are bounded with probability tending
to one. Inductively, these inequalities show that actual coefficients
converge to oracle coefficients, actual node errors vanish in normalized
\(L^2\), and actual norms remain bounded in probability. Local
Lipschitzness of scalar operations suffices by restricting to a compact
neighborhood of their deterministic limiting arguments. Equation
(\ref{eq-docs-special-data-limits-l3835}) transfers every oracle
empirical law to the actual program.


\bibliography{references.bib}



\end{document}
